% PDF 113--163；原书第 83--133 页。
\chapter{基于提示的学习}
\label{ch:prompting}
\chaptercredit{3}
\chapterabstract{本章讨论基于提示的学习，这是推动大语言模型近期进展的一项重要技术。开头将它与两种较早的方法作比较，即全监督学习（fully supervised learning）和预训练模型微调。接着细看基于提示的推理（prompt-based inference）需要哪些步骤，考察它如何用于从大语言模型提取知识，并概述它在自然语言处理中的应用。然后介绍提示工程的细节，讨论如何设计有效而高效的提示，将模型引向预期输出。此后转入答案工程，概述怎样优化模型回答，使之更准确、更切合上下文。还会讨论适用于较复杂查询的多提示技术。最后用两个教程说明基于提示的学习在实践中的优势。第一个教程演示，与传统的基于任务头的微调（head-based fine-tuning）相比，怎样用更少的训练样例取得更好的结果。第二个教程考察提示工程的不同方法。}

% PDF 113--120 上半页。
\section{引言}
语言模型的主要功能，是依据预训练过程中学到的语义表示，预测单个词元在其他词元组成的序列中出现的概率（第~\ref{ch:pretraining} 章）。在自回归任务中，可以利用这种能力生成语言。在掩码语言模型任务中，可以利用它做填空推理。两种方法都要先把一个初始词元序列送入模型进行推理。我们把这段起始序列称为提示。

\refcite{ch03-brown}{Brown 等人（2020）}的开创性研究表明，用提示就可以让大语言模型有效完成许多自然语言处理任务。研究者把一系列问答对连成长文本，最后留一个只有问题、没有答案的例子。图~\ref{fig:ch03-translation} 展示了一个英语译法语任务，其中使用相同的提示和不同数量的前置示例。除最后一项外，每个英语短语之后都跟着“\texttt{=>}”及法语译文。研究者用多个 GPT-3 变体，在应当填写答案的位置预测最可能出现的一个或多个词元。许多情况下，这些模型无需微调就能准确翻译句子。结果表明，通过提示可以利用大语言模型已经学到的语义知识来完成任务，无需分别针对每项任务再训练或调优。

\begin{figure}[!htbp]
\centering
% 重绘自原书第 84 页的图 3.1。提示和示例保留原文及拼写。
\begin{tikzpicture}[
  prompt/.style={fill=blue!10!gray!7,anchor=north west,inner sep=4pt,align=left,
    text width=53mm,font=\footnotesize\ttfamily},
  annotation/.style={anchor=west,font=\footnotesize\sffamily,inner sep=1.6pt},
  callout/.style={draw=black!61,line width=.48pt,-{Stealth[length=1.3mm]}}]
\node[prompt] at (0,0) {Translate English to French\\[5pt]
  cheese\hspace{8mm}=>};
\node[prompt] at (0,-1.78) {Translate English to French\\[5pt]
  sea otter\hspace{4mm}=> loutre de mar\\
  cheese\hspace{8mm}=>};
\node[prompt] at (0,-3.88) {Translate English to French\\[5pt]
  sea otter\hspace{4mm}=> loutre de mar\\
  peppermint\hspace{2mm}=> menthe poivrée\\
  plush girafe\hspace{1mm}=> girafe peluche\\
  cheese\hspace{8mm}=>};
\foreach \y/\note in {-.19/任务描述,-.87/提示,-1.97/任务描述,-2.63/示例,-3.02/提示,
  -4.07/任务描述,-5.15/示例,-5.99/提示} {
  \node[annotation] (tag) at (6.63,\y) {\note};
  \draw[callout] (tag.west) -- (5.47,\y);
}
\draw[black!61] (5.32,-4.69) -- (5.47,-4.69) -- (5.47,-5.70) -- (5.32,-5.70);
\end{tikzpicture}

\caption{通过提示进行英语到法语翻译的三个例子。每次都将示例和提示送入大语言模型，让模型预测接下来最可能出现的词。本例的答案为“fromage”。三个例子从上到下分别演示零样本、单样本和少样本推理。}
\label{fig:ch03-translation}
\end{figure}

提示的思想可以追溯到 \refcite{ch03-kumar}{Kumar 等人（2016）}提出的动态记忆网络（Dynamic Memory Network，DMN）。DMN 的神经网络架构用于处理输入序列和问题，建立情景记忆（episodic memory），生成相关答案（\refcite{ch03-xiong}{Xiong et al., 2016）}。与问题（提示）相对应的任务启动迭代注意力机制，使模型能够关注输入及前几轮迭代的结果。\refcite{ch03-radford-2019}{Radford 等人（2019）}揭示了这一思路的潜力。只要在足够大的数据集上训练，模型无需明确的监督就可以在多种自然语言处理任务上表现出较强能力。

此后出现了大量文献，研究基于提示的推理和学习的不同方法与改进。本章系统考察基于提示的推理的主要方面，包括基本流程、提示的形式、提示优化、答案空间设计，以及如何将它们用于不同的自然语言处理任务。\footnote{本章的术语和流程表述，主要参考 \refcite{ch03-liu-2023}{Liu 等人（2023）}的综述 \textit{Pre-train, Prompt, and Predict: A Systematic Survey of Prompting Methods in Natural Language Processing}。} 为交代提示方法的历史背景，先介绍过去几年影响较大的两种方法，即监督学习和预训练模型微调，并说明它们与基于提示的学习有何区别。

\subsection{全监督学习}
传统的做法是在有标签数据上训练自然语言处理模型。数据由输入与输出的配对样例组成，体现要学习的任务（\refcite{ch03-kotsiantis}{Kotsiantis et al., 2007）}。模型学习如何把输入映射为相应的输出，再从训练样例泛化，对未见过的数据进行预测。图~\ref{fig:ch03-supervised} 以逻辑回归分类器（logistic regression classifier）为例，模型学习句子和标签之间的关系。

\begin{figure}[!htbp]
\centering
% 重绘自原书第 86 页的图 3.2。示例文本和分类标签保留原样。
\begin{tikzpicture}[
  font=\footnotesize\sffamily,
  tag/.style={anchor=east,font=\footnotesize\sffamily\bfseries}]
\node[tag] at (1.92,3.98) {预测};
\node[tag] at (1.92,2.16) {模型};
\node[tag,align=right] at (1.92,.19) {训练\\数据};
\node[inner sep=0pt] (prediction) at (5.16,4.04) {
  \begin{tabular}{|l|}\hline
  Negative\\\hline Positive\\\hline
  \end{tabular}};
\filldraw[fill=black!4,draw=black!68] (3.55,1.02) rectangle (6.77,3.14);
\draw[black!50] (3.74,1.13) -- (3.74,3.00);
\draw[blue!38!gray,thick,domain=3.74:6.62,samples=60]
  plot (\x,{1.13+1.87/(1+exp(-5.8*(\x-5.40)))});
\node[font=\tiny,anchor=west] at (3.77,1.17) {0};
\node[font=\tiny,anchor=west] at (3.77,3.00) {1.0};
\node[align=left] at (4.77,2.15) {逻辑回归\\分类器};
\draw[black!62,line width=1.4pt,-{Stealth[length=2.1mm]}] (5.16,3.14) -- (prediction.south);
\node[inner sep=0pt,anchor=north] at (5.53,.91) {
  \begin{tabular}{|l|l|}\hline
  \textbf{文本} & \textbf{标签}\\\hline
  Cannot watch this movie. & Negative\\\hline
  Loved every bit of it. & Positive\\\hline
  \end{tabular}};
\end{tikzpicture}

\caption{根据训练数据做情感分类的逻辑回归分类器。通过最大化标签的似然，学习训练样例的特征权重。将学到的特征权重求和，再把和值送入逻辑函数，即 sigmoid 函数，得到介于 0 和 1 之间的概率。通过确定概率分布中的阈值，可将特定输入归为正面或负面。对于类别均衡的分类任务，常用阈值为 0.5。}
\label{fig:ch03-supervised}
\end{figure}

监督学习有明确的学习过程。有足够的有标签数据时，训练得到的模型在相应任务上通常有可靠的表现。因此，这种方法应用于情感分析（sentiment analysis）、机器翻译等多种自然语言处理任务。不过，它也有一些缺点。

首先，监督学习能否成功，很大程度上取决于有标签数据是否可获得、质量是否足够。这类数据可能稀缺，也可能需要高昂的标注成本或大量时间。其次，传统监督学习模型依赖专家进行特征工程（feature engineering），以定义用于预测的特征。这需要大量人工工作和丰富的专业知识，而且由于人们不完全了解特征在数据集中的自然分布，效率也受到限制（\refcite{ch03-och}{Och et al., 2004}；\refcite{ch03-zhang}{Zhang and Nivre, 2011}）。最后，监督学习模型较难泛化到训练数据范围之外，尤其遇到与训练集差异很大的样例时。

\begin{keypoint}
\noindent\textbf{优点}
\begin{itemize}
\item 表现可预测。
\item 适用面广。
\end{itemize}
\noindent\textbf{缺点}
\begin{itemize}
\item 高度依赖数据。
\item 需要特征工程。
\item 泛化能力有限。
\end{itemize}
\end{keypoint}

\subsection{预训练再微调}
这种方法先在大规模语料上无监督训练大语言模型，再用按目标任务标注的较小数据集作监督微调。模型由此适应特定任务，同时保有预训练获得的语义知识（\refcite{ch03-peters}{Peters et al., 2019}；\refcite{ch03-radford-2018}{Radford et al., 2018）}。图~\ref{fig:ch03-ptft} 以 BERT 模型为例，演示预训练再微调（Pre-Train and Fine-Tune，PTFT），BERT 见第~\ref{sec:ch02-bert} 节。

\sourcefigure{fig03_03.pdf}{以 BERT 说明预训练再微调。左侧（a）在涵盖丰富语言用法与主题的数据上作半监督预训练，学得丰富的语言语义表示。右侧（b）使用专为情感分类整理的数据作监督微调。预训练所学的语义知识得以迁移到微调任务，改善表现。}{fig:ch03-ptft}

这种方法在众多自然语言处理基准上达到过当时的先进水平，原因在于 PTFT 有几个主要优点。首先，预训练支持大量迁移学习，改善模型在不同自然语言处理任务上的表现（\refcite{ch03-kamath}{Kamath et al., 2019）}。其次，借助迁移学习，与全监督学习相比，微调对有标签数据的依赖较少，因此 PTFT 适用于资源有限的环境。最后，HuggingFace\footnote{\url{https://huggingface.co/}} 等平台和框架的发展，使自然语言处理领域的预训练模型微调流程越来越标准化，也越来越得到业界认可。

这些改进也带来一些代价。首先，训练和微调大型预训练模型需要大量算力，并非所有研究人员或开发者都能获得这些资源。其次，适合预训练的模型架构，如深度神经网络，往往难以解释，妨碍人们理解模型，并可能引起伦理问题。最后，预训练与微调的目标通常不同。预训练学习通用语义关系，微调则给文本分配标签。如果固定预训练模型，只微调新加的任务专用预测头，结果可能有所下降。

\begin{keypoint}
\noindent\textbf{优点}
\begin{itemize}
\item 可以利用迁移学习。
\item 表现优于全监督学习。
\item 较少依赖有标签数据。
\end{itemize}
\noindent\textbf{缺点}
\begin{itemize}
\item 算力需求大。
\item 语义空间与任务之间的映射不够精确。
\item 模型复杂，可解释性差。
\end{itemize}
\end{keypoint}

这种方法曾代表当时的先进水平，后来基于提示的学习成为调优大语言模型的新途径。

\subsection{基于提示的学习}
基于提示的学习又称基于提示的推理，是利用语言模型能力的一种新方法。通过提示，模型无需微调即可给出针对特定任务的回答。日常语境中，提示常被理解成人与模型沟通的一句话。例如，人向聊天机器人提出一个问题，这就是一条提示。提示也可以作为解决自然语言处理任务的策略。

\begin{figure}[!htbp]
\centering
% 重绘自原书第 88 页的图 3.4。提示、候选词和标签缩写取原图。
\begin{tikzpicture}[
  font=\footnotesize\sffamily,
  block/.style={fill=yellow!26,align=center,inner sep=2.2pt},
  flow/.style={draw=black!62,line width=.51pt,-{Stealth[length=1.5mm]}}]
\node[text width=28mm,align=left,inner sep=0pt,font=\scriptsize\ttfamily] (prompt) at (0,0)
  {Cannot watch this movie. This is [MASK]};
\node[block] (model) at (3.14,0) {
  \begin{tabular}{@{}c@{}}\textbf{语言模型}\\[2pt]
  \begin{tabular}{|lc|}\hline
  \rowcolor{blue!13}\textbf{词元} & $P$\\\hline
  \rowcolor{blue!7}terrible & 0.97\\\hline
  \rowcolor{blue!7}great & 0.03\\\hline
  \end{tabular}
  \end{tabular}};
\node[block] (vocab) at (6.13,0) {
  \begin{tabular}{@{}c@{}}\textbf{词表}\\[2pt]
  \begin{tabular}{|lc|}\hline
  \rowcolor{blue!13}\textbf{词元} & \textbf{类别}\\\hline
  \rowcolor{blue!7}terrible & neg\\\hline
  \rowcolor{blue!7}great & pos\\\hline
  \end{tabular}
  \end{tabular}};
\node[align=left] (result) at (8.84,0) {负面\\类别};
\draw[flow] (prompt.east) -- (model.west);
\draw[flow] (model.east) -- (vocab.west);
\draw[flow] (vocab.east) -- (result.west);
\end{tikzpicture}

\caption{用“预训练、提示、预测”做分类。将一条留有空缺（\texttt{[MASK]}）的提示送入预训练语言模型，观察模型预测某些词填入空缺的概率。取概率最高的词元“terrible”，返回它所映射的标签类别“negative”（负面）。}
\label{fig:ch03-mask}
\end{figure}

在这类应用中，提示是一段自然语言文本，留出一个或多个词的位置，由大语言模型依据内部的语义模型补全。图~\ref{fig:ch03-mask} 给出了示例。无需微调模型来直接预测正面或负面标签，可以把句子“\texttt{Cannot watch this movie. This is [MASK].}”交给大语言模型，再看“great”与“terrible”哪一个更可能填入掩码位置。这里模型会认为“terrible”更可能，因为这样补成的句子比填入一个正面词更符合语义。这种预测可以再映射到标签类别，本例为“negative”。

基于提示的预测也可以做开放式生成。例如，使用“\texttt{English: Cannot watch this movie. German: }”作为翻译提示，让模型继续预测末尾的词元，模型就会生成输入句子的德语译文。

这种方法需要组织提示，引导语言模型产生指定自然语言处理任务所需的输出。它利用预训练语言模型生成连贯文本的能力，减少所需的有标签数据，并支持零样本和少样本学习。因此，基于提示的学习成为活跃的研究领域，在多项自然语言处理任务中表现良好。

\begin{practicaltip}
本书用零样本、单样本和少样本描述基于提示的学习所用训练数据集的不同规模。在零样本设置中，预训练模型无须使用额外的训练样例就可以完成目标任务。在单样本与少样本设置中，用一个特定于任务的训练样例，或少量这样的样例（少样本，$\lesssim100$）引导模型。与可能需要大量样本的其他技术相比，基于提示的模型用有限训练数据便能取得不错的表现，这是一个重要优点。
\end{practicaltip}

基于提示的学习有几项优点。只要能把任务写成填空问题，就能适配多种自然语言处理任务。与 PTFT 相比，预测任务与预训练学到的语义知识衔接得更紧密。由于模型权重没有更新，适应新信息时，模型不会“忘记”旧信息。由此，同一组原始权重可用于多个场景，大幅降低训练成本，也简化了应用部署。最后，利用语言模型的文本生成能力，可以减少对有标签数据的需求，从而支持零样本或少样本学习。

另一方面，提示的选择可能很大程度上影响模型表现，因此提示工程十分重要，也可能颇具挑战。与预训练模型微调一样，基于提示的学习背后的机制也可能难以理解和解释。最后，基于提示的模型可能给出不一致的输出，尤其是在没有充分接触目标任务或提示含义不明确时。

\begin{keypoint}
\noindent\textbf{优点}
\begin{itemize}
\item 适用于多种自然语言处理任务。
\item 训练和推理的语义衔接较好。
\item 所需的有标签数据较少。
\end{itemize}
\noindent\textbf{缺点}
\begin{itemize}
\item 对提示设计很敏感。
\item 可解释性有限。
\item 结果可能不一致。
\end{itemize}
\end{keypoint}

后续各节介绍基于提示的学习的基本概念，说明怎样构造有效的提示，并讨论如何用不同的提示方法微调大语言模型。

% PDF 120--127；表 3.1 见 PDF 128--129。
\section{基于提示的学习基础}
\subsection{基于提示的学习的形式化描述}
对语言模型使用提示，就是给模型一个初始输入或一系列文本线索，模型再根据这些输入生成符合上下文且连贯的回答。以情感分析中的影评句子“Cannot watch this movie”为例。本节用这句话比较监督学习、预训练再微调和基于提示的学习，说明它们的不同之处。

如图~\ref{fig:ch03-formal-supervised} 所示，监督学习训练参数为 $\theta$ 的模型，使其从训练集中的配对样例 $(\symbf{x},\symbf{y})$ 学习。这个例子中，与 $(\symbf{x},\symbf{y})$ 对应的训练样例是（Cannot watch this movie，negative）。

\begin{figure}[!htbp]
\centering
% 重绘自原书第 90 页的图 3.5。
\begin{tikzpicture}[font=\small\sffamily,
  flow/.style={draw=black!60,-{Stealth[length=1.4mm]},line width=.48pt},
  bracket/.style={draw=black!49,line width=.45pt}]
\node[font=\footnotesize\itshape] (input) at (0,0) {Cannot watch this movie.};
\node[fill=black!15,align=center,minimum width=21mm] (classifier) at (4.00,0)
  {分类器\\$P(\symbf{y}\mid\symbf{x};\theta)$};
\node[font=\footnotesize\itshape] (output) at (6.51,0) {negative};
\draw[flow] (input.east) -- (classifier.west);
\draw[flow] (classifier.east) -- (output.west);
\draw[bracket] (input.south west) -- ++(0,-.19) -| (input.south east);
\draw[bracket] (output.south west) -- ++(0,-.19) -| (output.south east);
\node[below=2mm of input] {$\symbf{x}$};
\node[below=2mm of output] {$\symbf{y}$};
\end{tikzpicture}

\caption{监督学习表示为给定输入 $\symbf{x}$ 时输出 $\symbf{y}$ 的概率。}
\label{fig:ch03-formal-supervised}
\end{figure}

在预训练再微调中，修改输入以包含一个头词元 \texttt{[CLS]}。它对应模型试图预测的类别标签，如“positive”或“negative”，分别对应正面和负面的影评。微调也称基于任务头的调优（head-based tuning），用标签和输入学习这个分类器的参数 $\theta$，见图~\ref{fig:ch03-formal-ptft}。

\begin{figure}[!htbp]
\centering
\scalebox{.99}{% 重绘自原书第 91 页的图 3.6。
\begin{tikzpicture}[
  x=.97cm,font=\footnotesize\sffamily,
  block/.style={fill=black!14,align=center,inner sep=2pt,minimum width=14mm,minimum height=8mm},
  flow/.style={draw=black!60,-{Stealth[length=1.3mm]},line width=.48pt},
  bracket/.style={draw=black!49,line width=.45pt}]
\node[font=\footnotesize\itshape] (input) at (0,0) {Cannot watch this movie.};
\node[block] (ptm) at (3.22,0) {预训练\\模型};
\node[block] (head) at (4.82,0) {微调\\分类器};
\node[draw=blue!28!gray!20,fit=(ptm)(head),inner sep=1.8pt] (model) {};
\node[anchor=west,font=\footnotesize\itshape,inner sep=0pt] (output) at (6.10,0)
  {[CLS] Cannot watch this movie.};
\node[anchor=west,font=\footnotesize\itshape,inner sep=0pt] at (6.10,.46) {negative};
\draw[flow] (input.east) -- (model.west);
\draw[flow] (model.east) -- (output.west);
\draw[bracket] (input.south west) -- ++(0,-.20) -| (input.south east);
\draw[bracket] (6.10,-.17) -- (6.10,-.35) -- (6.82,-.35) -- (6.82,-.17);
\node at (0,-.58) {$\symbf{x}$};
\node at (6.46,-.58) {$\symbf{y}$};
\end{tikzpicture}
}
\caption{预训练模型加微调。分类头作用于预训练模型对输入 $\symbf{x}$ 的嵌入，产生输出 $\symbf{y}$。}
\label{fig:ch03-formal-ptft}
\end{figure}

基于提示的学习则把每个输入放入一个提示模板（prompt template）。模板既容纳输入，也为预测输出留出与推理任务相关的空位。例如，图~\ref{fig:ch03-formal-prompt} 的例子在输入之后附加“It is [z]”，再根据语言模型的概率，找出最可能填入空位的词元 $z$。

因此，基于提示的推理包含两个主要部分。一是把下游任务转换为语言建模问题的模板，二是把语言模型文本输出转换成分类标签的一组标签词（label word）。这里选“great”和“terrible”作为两个可能输出。与“positive”和“negative”相比，这两个词填入提示更加自然，也更有可能被模型预测出来。图~\ref{fig:ch03-formal-prompt} 展示了模板变换和标签词的选择。这种做法无须引入新参数。

\begin{figure}[!htbp]
\centering
% 重绘自原书第 91 页的图 3.7。
\begin{tikzpicture}[
  font=\footnotesize\sffamily,
  block/.style={fill=black!14,align=center,inner sep=2pt,minimum width=13mm,minimum height=8mm},
  flow/.style={draw=black!60,-{Stealth[length=1.3mm]},line width=.48pt},
  bracket/.style={draw=black!49,line width=.45pt}]
\node[font=\footnotesize\itshape] (input) at (0,0) {Cannot watch this movie.};
\node[block] (lm) at (2.90,0) {语言\\模型};
\node[block] (prompter) at (4.36,0) {提示\\流程};
\node[draw=blue!28!gray!20,fit=(lm)(prompter),inner sep=1.8pt] (model) {};
\node[anchor=west,text width=38mm,align=left,font=\footnotesize\itshape,inner sep=0pt] (output) at (5.45,0)
  {Cannot watch this movie. This is [z].};
\draw[flow] (input.east) -- (model.west);
\draw[bracket] (input.south west) -- ++(0,-.16) -| (input.south east);
\draw[bracket] (output.south west) -- ++(0,-.16) -| (output.south east);
\node[below=2mm of input] {$\symbf{x}$};
\node[below=2mm of output] {$\symbf{x}'$};
\node[font=\footnotesize\rmfamily] at (3.63,-.89)
  {$[z]\in\{\text{“great”},\text{“terrible”}\}$};
\end{tikzpicture}

\caption{基于提示的学习。用提示函数将任务由 $\symbf{x}$ 改写成 $\symbf{x}'$，使语言模型输出能够映射到 $\symbf{y}$。}
\label{fig:ch03-formal-prompt}
\end{figure}

形式上，若标签词映射 $\mathcal{M}$ 将标签 $\symbf{y}$ 映射到单个词，给定原始输入 $\symbf{x}$ 与加提示后的输入 $\symbf{x}'$，可将基于提示的微调过程写为
\begin{equation}
p(\symbf{y}\mid\symbf{x})=p([\symbf{z}=\mathcal{M}(\symbf{y})\mid\symbf{x}']
\tag{3.1}
\end{equation}
\begin{equation}
p(\symbf{y}\mid\symbf{x})=
\frac{\exp(\symbf{w}_{\mathcal{M}(\symbf{y})}\cdot\symbf{h}_z)}
{\sum_{\symbf{y}'\in\symbf{y}}\exp(\symbf{w}_{\mathcal{M}(\symbf{y}')}\cdot\symbf{h}_z)}.
\tag{3.2}
\end{equation}
这里 $\symbf{h}_z$ 是答案 $z$ 的隐藏向量，$\symbf{w}$ 是与映射到 $\symbf{y}$ 的词子集有关的 softmax 之前的向量。
\transnote{式（3.1）末尾缺少与 $p$ 后开括号对应的右圆括号。式（3.2）的求和域按原书第 91 页保留。}

\begin{practicaltip}
基于提示的方法往往比第 2 章讨论的基于任务头的微调表现更好，尤其是在数据较少时。原因是 BERT 等预训练模型加入了随机初始化的新参数，只用有限的样例难以有效优化这些参数。第 3.6 节的教程将给出具体例子。
\end{practicaltip}

\subsection{基于提示的学习流程}
仍以影评情感分类为例，说明数据流和输出类别预测的各个步骤。如图~\ref{fig:ch03-flow} 所示，基于提示的学习有三个主要步骤。
\begin{enumerate}
\item 提示函数（prompting function）$f_{\mathrm{prompt}}(\cdot)$ 将输入文本 $\symbf{x}$ 改写为提示 $\symbf{x}'(z)=f_{\mathrm{prompt}}(\symbf{x})(z)$。
\item 函数 $f_{\mathrm{fill}}(\symbf{x}';z)$ 找出最可能填入空位 $z$ 的词元。
\item 用映射函数把得分最高的答案 $z^*$ 与得分最高的输出 $\symbf{y}$ 对应起来。
\end{enumerate}
在这个例子中，输入句子“Cannot watch this movie”记为 $\symbf{x}$。模板是“[x] It was a [z] movie”，输出是情感预测“negative”，记为 $\symbf{y}$。下面分别讨论这些步骤。

\begin{figure}[!htbp]
\centering
% 重绘自原书第 93 页的图 3.8。
\begin{tikzpicture}[
  font=\small\sffamily,
  heading/.style={fill=black!18,inner sep=2.5pt,font=\small\sffamily\bfseries},
  process/.style={anchor=west,align=left,inner sep=1pt},
  flow/.style={line width=.66pt,draw=black!65,-{Stealth[length=1.75mm]}},
  data/.style={font=\small\rmfamily\itshape,align=center,inner sep=2pt}]
\node[heading] at (1.10,5.69) {提示流程};
\node[heading] at (6.53,5.69) {数据流};
\draw[flow,densely dotted] (2.84,5.78) -- (2.84,-.22);
\node[process] at (0,5.03) {输入 $\symbf{x}$};
\node[process] at (0,3.98) {模板};
\node[process] at (0,2.86) {提示函数\\$\symbf{x}'=f_{\mathrm{prompt}}(\symbf{x})$};
\node[process] at (0,1.74) {预测\\$f_{\mathrm{filled}}(\symbf{x}',z^*)$};
\node[process] at (0,.40) {答案映射\\$z^*\Rightarrow \symbf{y}$};
\node[data] (input) at (6.53,5.03) {Cannot watch this movie.};
\node[data] (template) at (6.53,3.98) {[x] It was a [z] movie.};
\node[data] (prompt) at (6.53,2.86) {Cannot watch this movie. It was a [z] movie.};
\node[data] (filled) at (6.53,1.74) {Cannot watch this movie. It was a \colorbox{black!15}{terrible} movie.};
\node[data] (class) at (6.53,.40) {\texttt{terrible => Negative}};
\foreach \origin/\target in {input/template,template/prompt,prompt/filled,filled/class} {
  \draw[flow] (\origin.south) -- (\target.north);
}
\end{tikzpicture}

\caption{基于提示的学习流程。先将提示函数应用于输入，再填补所得文本中的掩码位置，并将输出映射到相应的答案。}
\label{fig:ch03-flow}
\end{figure}

\subsubsection{添加提示}
添加提示（prompt addition）可以表示成函数 $f_{\mathrm{prompt}}(\cdot)$。它接收输入文本 $\symbf{x}$ 和带答案空位 $z$ 的指定模板，生成提示 $\symbf{x}'=f_{\mathrm{prompt}}(\symbf{x})$。该函数包含两个阶段。
\begin{enumerate}
\item 使用预先确定的模板。模板是一段含两个空位的文本，输入空位 \texttt{[x]} 用于放入输入，答案空位 \texttt{[z]} 用于放入生成的中间答案文本。之后再把这个中间文本映射为输出 $\symbf{y}$。
\item 用给定输入替换输入空位 \texttt{[x]}。
\end{enumerate}
沿用影评情感分类的例子，用模板“[x] It was a [z] movie”转换输入句子“Cannot watch this movie.”，就会得到提示“Cannot watch this movie. It was a [z] movie.”。

\subsubsection{答案搜索}
下一步计算提示中指定空位所有可能答案的概率。这个过程常称为答案搜索（answer search），是用模型进行推理的阶段。对上述提示模板，可以分别用“great”和“terrible”作为正面和负面类别的候选答案。概率最高的填空词元就是预测答案。

形式上，答案搜索主要分两步。
\begin{enumerate}
\item 函数 $f_{\mathrm{fill}}(\symbf{x}',z)$ 把一个候选答案 \texttt{[z]} 填进指定空位。根据任务，可以从模型的整个词表或一个较小的词子集 $Z$ 中选取答案。这个过程得到的文本称为已填充提示（filled prompt）。
\item 使用搜索函数，如 $\arg\max$ 搜索，找出得分最高的输出。做法是用预训练语言模型 $P(\mathord{\,};\theta)$ 计算相应已填充提示的概率。搜索函数选择概率最高的输出，找出最符合上下文、语义最连贯的答案，以满足提示要求。
\end{enumerate}
这个搜索函数可写为
\begin{equation}
\hat{z}=\operatorname*{search}_{z\in Z} P(f_{\mathrm{fill}}(\symbf{x},z);\theta).
\tag{3.3}
\end{equation}
这里给定模板函数 $f_{\mathrm{fill}}$ 和模型参数 $\theta$，在所有可能的答案 $Z$ 中找出概率最高的答案 $z$。然后可把输出 $\hat{z}$ 映射到更自然、更容易解释的答案空间。
\transnote{式（3.3）按原式使用 $\symbf{x}$，前文答案搜索的定义则使用 $\symbf{x}'$。}

\subsubsection{答案映射}
在答案搜索中得到中间文本或词元后，需要把它们映射为最终所需的输出形式。这称为答案映射（answer mapping）。

例如，在影评情感分析任务中，基于提示的模型可能生成“terrible”或“great”等中间文本来填空。答案映射随后把这些文本对应到为任务定义的情感标签或数值，例如用 1 表示正面，0 表示负面。形式上，映射把得分最高的答案 $\hat{z}$ 变为得分最高的输出 $\hat{\symbf{y}}$。

这一步之所以必要，是因为最适合填入模板的自然用词，未必与标签完全一致。例如，餐厅评论的情感分类通常用一星到五星表示。“one star”作为答案没有“terrible”自然，因此提示中可以使用后者作为候选，之后再把它映射到“一星”类别。答案搜索所用的词元，必须转换成符合任务目标的形式。

% PDF 124--127；表 3.1 见 PDF 128--129。
\subsection{基于提示的知识提取}
语言模型在预训练中既学习一般性的语义知识，也学到无数具体的事实性陈述，这些知识通常称为参数化知识（parametric knowledge）。图~\ref{fig:ch03-formal-prompt} 的提示任务只需了解词语之间的语义关系即可完成。再考虑另一个问题，“What city was Dante Alighieri born in?”，也就是问但丁·阿利吉耶里（Dante Alighieri）出生在哪座城市。仅凭语义无法回答这个问题，模型必须接触过具体的事实信息。大语言模型预训练语料库的巨大规模使这一点成为可能。例如，如果语料库含有但丁的维基百科词条，模型权重就会受到这样的句子影响。“但丁出生在佛罗伦萨共和国的佛罗伦萨，那里现在位于意大利。”恰当设计的提示可以让模型复述这类参数化知识。例如，
\begin{equation}
\text{\ttfamily Dante was born in [MASK]}
\tag{3.4}
\end{equation}
由于模型训练时接触过能回答这个问题的数据，就会认为“Florence”很可能填在掩码词元的位置。这说明了大语言模型基于提示的推理的一项基本潜力，即可以将大语言模型作为知识库（knowledge base）使用。

构建传统知识库需要从不同数据源提取关系知识，也需要开发实体抽取、共指消解（coreference resolution）、实体链接（entity linking）和关系抽取（relation extraction）等自然语言处理流程（\refcite{ch03-petroni}{Petroni et al., 2019）}。各环节都有难点，错误难以避免。错误沿流程传播和积累，可能显著影响所建知识库的实用性（\refcite{ch03-petroni}{Petroni et al., 2019）}。图~\ref{fig:ch03-kb} 从概念上比较了两种做法。

\begin{figure}[!htbp]
\centering
\scalebox{.99}{% 重绘自原书第 95 页的图 3.9。
\begin{tikzpicture}[
  x=.97cm,font=\footnotesize\sffamily,
  head/.style={font=\footnotesize\sffamily\bfseries},
  vertex/.style={circle,inner sep=0pt,minimum size=3.2mm,draw=black!76,line width=.7pt},
  wire/.style={draw=black!68,line width=.55pt},
  flow/.style={wire,-{Stealth[length=1.4mm]}},
  memory/.style={fill=yellow!22,align=center,minimum width=23mm,minimum height=8mm,inner sep=2pt}]
\node[head] at (1.95,4.96) {记忆};
\node[head] at (5.97,4.96) {查询};
\node[head] at (9.80,4.96) {答案};
\node[head,anchor=east] at (.05,3.47) {知识库};
\node[vertex,fill=blue!35] (base) at (.40,3.47) {};
\node[vertex,fill=green!30] (top) at (.95,4.12) {};
\node[vertex,fill=green!30] (dante) at (1.85,3.47) {};
\node[vertex,fill=green!30] (bottom) at (1.19,2.83) {};
\node[vertex,fill=yellow!63] (topentity) at (2.65,4.20) {};
\node[vertex,fill=yellow!63] (florence) at (2.65,2.66) {};
\draw[wire] (base)--(top)--(bottom)--(base)--(dante)--(topentity);
\draw[wire] (dante) -- node[font=\tiny\ttfamily,fill=white,inner sep=.8pt,sloped] {born-in} (florence);
\node[anchor=west,font=\scriptsize] at (2.04,3.49) {Dante};
\node[anchor=west,font=\scriptsize] at (2.83,2.65) {Florence};

\node[memory] (kb) at (5.97,3.41) {访问符号记忆};
\node[font=\scriptsize\ttfamily] (rel) at (5.97,4.23) {(Dante, born-in, X)};
\node at (9.80,3.41) {Florence};
\draw[wire] (3.53,3.41) -- (kb.west);
\draw[flow] (rel.south) -- (kb.north);
\draw[flow] (kb.east) -- (9.11,3.41);

\draw[draw=black!62,densely dotted,line width=1pt] (-.62,2.13) -- (10.48,2.13);
\node[head,anchor=east] at (.05,.88) {语言模型};
\foreach \offset in {.32,.16,0} {
  \begin{scope}[shift={(\offset,\offset)}]
  \filldraw[draw=black!25,fill=black!6] (.58,.35) rectangle (1.82,1.59);
  \foreach \x/\y/\tone in {.82/.59/blue,.82/.97/green,.82/1.35/green,
   1.20/.59/green,1.20/.97/yellow,1.20/1.35/blue,
   1.58/.59/green,1.58/.97/yellow,1.58/1.35/yellow} {
    \filldraw[draw=black!67,fill=\tone!28] (\x,\y) circle[radius=1.88mm];
  }
  \end{scope}
}
\node[font=\scriptsize] at (1.28,.03) {如 ELMo/BERT};
\node[memory] (lm) at (5.97,.80) {访问神经语言\\模型记忆};
\node[font=\scriptsize\ttfamily] (prompt) at (5.97,1.68) {"Dante was born in [MASK]"};
\node at (9.80,.80) {Florence};
\draw[wire] (3.53,.80) -- (lm.west);
\foreach \x/\y in {4.96/5.28,5.55/5.69,6.11/6.11,6.57/6.46,7.08/6.83} {
  \draw[flow] (\x,1.49) -- (\y,1.21);
}
\draw[flow] (lm.east) -- (9.11,.80);
\end{tikzpicture}
}
\caption{向传统知识库和语言模型查询事实信息。虚线上方是为查询实体关系而专门设计的知识库，本例关系是“但丁出生在佛罗伦萨”。下方的语言模型原本用于根据上下文预测被遮盖的词。因此，可借助提示让它报告预训练时接触过的事实。}
\label{fig:ch03-kb}
\end{figure}

\begin{keypoint}
与使用传统知识库相比，给大语言模型提供提示有几个主要优点。首先，它不需要预先定义模式（schema-free），关系知识是预训练的涌现结果，存于语言模型内部，无须专为这项任务开发模型。其次，现代预训练语料覆盖的信息范围很广，模型的泛化能力也很强。理论上，同一个语言模型可以支持多种知识应用和范围更广的常见自然语言处理任务。
\end{keypoint}

\begin{practicaltip}
基于大语言模型提取知识也面临很大的挑战和风险。要解释模型回答并不容易，因为往往无法追查是哪些因果过程使基于提示的推理产生某个特定回答。回答的准确性也需要验证。我们试图引出的知识，是预训练过程中涌现的，模型并未接受专门训练以学习这些事实。预训练数据的规模也大到难以验证知识或按质量排序。若数据取自互联网，同样缺乏对知识来源和可靠性的审查，可能使模型学到不真实的“事实”。因此，用户应保持合理的怀疑态度，采用适当的评估方法防范这些错误。最后，基于提示的知识库能否给出一致结果，很大程度上取决于提示设计的质量。第 3.3 节将详细讨论提示优化的方法。
\end{practicaltip}

\subsection{基于提示的学习在各类自然语言处理任务中的应用}
许多自然语言处理任务都适合用提示来完成。本节列出常见任务，并介绍它们的输入、模板、提示及答案映射。理解这些要素以及它们之间的关系，有助于了解如何把各项自然语言处理任务有效转化成基于提示的学习问题并加以求解。

我们将这些任务分成三类。
\begin{enumerate}
\item \textbf{文本分类。}为给定的输入句子分配适当的类别标签。提示接收输入句子，留出专门位置供模型生成中间答案，随后可把中间答案映射为分类标签。
\item \textbf{标注（tagging）。}为给定文本中的个别元素，如词或短语，分配标签或标记。提示先包含待标注元素所在的文本，再就该元素具体提问，给出若干选项供模型选择。
\item \textbf{文本生成。}根据提示中的任务生成一段文本，通常包含不止一个词元。提示给出相关上下文，如待摘要的段落或待翻译的句子，再用明确指令告诉模型如何处理这段上下文。
\end{enumerate}

表~\ref{tab:ch03-tasks} 列出三类共 17 项任务，简述各项任务，并给出可以用于完成该任务的示例输入、模板和答案空间。多样的用例体现了基于提示的学习的灵活性。不过，提示需要根据每项任务精心设计。下一节将把提示设计拆分为几个可以分别优化的方面，以改善基于提示的学习的表现。

% 按原书第 98--99 页重排表 3.1。
\begingroup
\small
\setlength{\tabcolsep}{4pt}
\setlength{\LTcapwidth}{\linewidth}
\renewcommand{\arraystretch}{1.07}
\newcommand{\taskexample}[3]{%
  \textbf{输入}\enspace #1\par
  \textbf{模板}\enspace #2\par
  \textbf{答案}\enspace #3}
\newcommand{\inputslots}{[$\symbf{x}_1,\ldots,\symbf{x}_n$]}
\begin{longtable}{@{}>{\raggedright\arraybackslash}p{.37\linewidth}>{\raggedright\arraybackslash}p{.60\linewidth}@{}}
\caption{基于提示的自然语言处理方法。每行左侧给出任务及定义，右侧给出示例，其中包括待处理的输入句子、用于基于提示的推理的模板，以及一种可能的答案空间。任务分为文本分类、标注与文本生成三类。}
\label{tab:ch03-tasks}\\
\endfirsthead
\multicolumn{2}{@{}l}{\textbf{表~\ref{tab:ch03-tasks}（续）}}\\
\toprule
\textbf{任务} & \textbf{示例}\\
\midrule
\endhead
\midrule
\multicolumn{2}{r@{}}{\textit{续下页}}\\
\endfoot
\bottomrule
\endlastfoot
\multicolumn{2}{@{}l}{\textbf{文本分类}}\\
\toprule
\textbf{任务} & \textbf{示例}\\
\midrule
\textbf{情感分析}\enspace 将文本的情感判为正面、负面或中性。 &
\taskexample{I hate this movie.}{\texttt{[x]} It was a \texttt{[z]} movie.}{great, terrible, \ldots}\\
\midrule
\textbf{作者归属识别}\enspace 从预先给定的作者集合中判断给定文本的作者。 &
\taskexample{It was the best of times, it was the worst of times, it was the age of wisdom, it was the age of foolishness}
{The author of \texttt{[x]} is most likely \texttt{[z]}.}{Dickens, Carroll, Austin, \ldots}\\
\midrule
\textbf{垃圾消息检测}\enspace 判断邮件或短信是否属于垃圾消息。 &
\taskexample{Congratulations! You have won! Click here to claim your free vacation.}
{This message: \texttt{[x]} is classified as \texttt{[z]}.}{Spam, Non-Spam}\\
\midrule
\textbf{情绪分类}\enspace 从预定义的情绪集合中，判断文本表达了哪种情绪。 &
\taskexample{I just won the lottery!}{This text: \texttt{[x]} expresses the emotion \texttt{[z]}.}
{anger, surprise, sadness, happiness}\\
\midrule
\textbf{意图检测}\enspace 判断用户发出查询或消息的意图，常用于聊天机器人和虚拟助手。 &
\taskexample{What’s the weather like today?}{\texttt{[x]} The user’s intent is \texttt{[z]}.}
{\texttt{get\_weather}, \texttt{set\_alarm}}\\
\midrule
\textbf{语种识别}\enspace 判断给定文本以何种语言写成。 &
\taskexample{¿Cómo estás?}{\texttt{[x]} The language is \texttt{[z]}.}{Spanish, French, \ldots}\\
\midrule
\textbf{仇恨言论检测}\enspace 判断给定文本是否含有仇恨言论。 &
\taskexample{I can’t stand them.}{\texttt{[x]} The text contains \texttt{[z]} speech.}{hate, non-hate}\\
\addlinespace[5pt]
\multicolumn{2}{@{}l}{\textbf{标注}}\\
\toprule
\textbf{任务} & \textbf{示例}\\
\midrule
\textbf{词性标注（POS）}\enspace 给词分配语法类别，如名词、动词、形容词和副词。 &
\taskexample{She is running in the park.}{In the sentence \inputslots, the word [$\symbf{x}_i$] has POS-tag [$\symbf{z}_j$].}{noun, verb, adjective, \ldots}\\
\midrule
\textbf{命名实体识别（NER）}\enspace 识别并分类文本中提到的人物、日期、地点、机构等实体。 &
\taskexample{John met Mary in London.}{In the sentence \inputslots, the word [$\symbf{x}_i$] the named entity label is [$\symbf{z}_j$].}{location, organization, \ldots}\\
\midrule
\textbf{组块分析或浅层句法分析}\enspace 依据语法结构把相邻的词或词元组成较大的单元，即组块，例如名词短语或动词短语。 &
\taskexample{She is running in the park.}{In the sentence \inputslots, the word [$\symbf{x}_i$] the chunk label is [$\symbf{z}_j$].}{’B-VP’（动词短语起点），’I-VP’（动词短语内部），\ldots}\\
\midrule
\textbf{依存句法分析}\enspace 找出句子中词语之间的句法依存关系，包括标注主语、宾语、修饰语等，并显示它们之间的关系。 &
\taskexample{She is running in the park.}{In the sentence \inputslots, the word [$\symbf{x}_i$] the dependency relation is [$\symbf{z}_j$].}{’nsubj’（名词性主语），’root’（句子的根），’dobj’（直接宾语），\ldots}\\
\midrule
\textbf{成分句法分析或短语结构分析}\enspace 识别句子的成分结构，即词如何组成名词短语（NP）、动词短语（VP）等语法短语。 &
\taskexample{She is running in the park.}{In the sentence \inputslots, the word [$\symbf{x}_i$] the constituent category is [$\symbf{z}_j$].}{’NP’（名词短语），’PP’（介词短语），’VP’（动词短语），\ldots}\\
\midrule
\textbf{语义角色标注（SRL）}\enspace 按照词或短语与谓词（通常为动词）之间的语义关系，为它们分配施事、工具等角色。 &
\taskexample{John gave Mary a book.}{In the sentence \inputslots, the word [$\symbf{x}_i$] has semantic role [$\symbf{z}_j$].}{Agent, Theme, Location, \ldots}\\
\midrule
\textbf{共指消解}\enspace 识别文本中指向同一实体的词语或短语，并将它们关联起来。 &
\taskexample{Jane is a talented software engineer. She was recently promoted to team lead.}{In the text with words: \inputslots, does the word [$\symbf{x}_i$] refer to the word $\symbf{x}_j$?}{Yes, No}\\
\addlinespace[5pt]
\multicolumn{2}{@{}l}{\textbf{文本生成}}\\
\toprule
\textbf{任务} & \textbf{示例}\\
\midrule
\textbf{摘要生成}\enspace 给定一篇长文本，生成保留原文要点或关键信息的较短版本。 &
\taskexample{<Long text to be summarized.>}{Please provide a summary for the following text: \texttt{[x]}. Summary: \texttt{[z]}.}{<summarized version of the long text>}\\
\midrule
\textbf{问答}\enspace 给定问题和上下文，依据上下文中的信息生成答案。 &
\taskexample{<Context or passage>, <question>}{Here is the context: \texttt{[x]} What is the answer to the question: \texttt{[w]} ? Answer: \texttt{[z]} .}{<answer to the question based on the context>}\\
\midrule
\textbf{机器翻译}\enspace 将文本从一种语言译成另一种语言，保留原意和语境。 &
\taskexample{¿Cómo estás?}{Translate the following text from the source language to the target language: \texttt{[x]} Translation: \texttt{[z]}.}{<translated text in target language>}\\
\end{longtable}
\endgroup


% PDF 127、130--133；原书第 97、100--103 页。
\section{提示工程}
上一节用几个简单的例子说明，如何利用提示解决不同的自然语言处理任务。提示的具体措辞对取得良好结果很重要。为提高模型在目标下游任务上的表现而设计合适的提示函数，就是提示工程。设计提示时，需要考虑并综合多个要素，包括选择预训练模型、确定合适的提示形式、设计提示模板和答案工程。模板的设计大致分为两类。
\begin{itemize}
\item 手工设计模板。
\item 自动设计模板。
\end{itemize}
前者借助人的专业知识和反复尝试来优化提示，后者用各种自动化方法寻找适合给定任务的模板。图~\ref{fig:ch03-design} 概览了接下来两个节的结构。下面介绍提示分类的基本术语，概述手工提示工程，并详细讨论文献中的几种自动化方法。

\begin{figure}[!htbp]
\centering
% 重绘自原书第 100 页的图 3.10。
\begin{tikzpicture}[
  font=\footnotesize\sffamily,
  option/.style={align=center,minimum width=24mm,minimum height=8mm,inner sep=2pt},
  wire/.style={draw=black!48,line width=.65pt}]
\node[option,fill=green!17] (root) at (0,0) {基于提示\\的学习};
\node[option,fill=yellow!65!orange] (model) at (2.92,3.80) {预训练模型};
\node[option,fill=blue!15] (eng) at (2.92,2.70) {提示工程};
\node[option,fill=orange!20] (ans) at (2.92,-2.00) {答案工程};
\node[option,fill=black!10] (multi) at (2.92,-3.10) {多提示学习};
\node[option,fill=black!55,text=white] (tuning) at (2.92,-4.20) {基于提示的微调};

\node[option,fill=blue!26] (promptshape) at (5.84,3.80) {提示形式};
\node[option,fill=blue!26] (template) at (5.84,1.60) {提示模板设计};
\node[option,fill=orange!30] (answershape) at (5.84,-.60) {答案形式};
\node[option,fill=orange!30] (search) at (5.84,-3.30) {答案搜索};

\node[option,fill=blue!37] (cloze) at (8.76,4.35) {完形填空提示};
\node[option,fill=blue!37] (prefix) at (8.76,3.25) {前缀提示};
\node[option,fill=blue!59!black,text=white] (manual) at (8.76,2.30) {手工设计};
\node[option,fill=blue!59!black,text=white] (auto) at (8.76,1.38) {自动设计};
\node[option,fill=orange!59] (token) at (8.76,.50) {词元};
\node[option,fill=orange!59] (span) at (8.76,-.60) {片段};
\node[option,fill=orange!59] (sentence) at (8.76,-1.70) {句子};
\node[option,fill=orange!68!black,text=white] (human) at (8.76,-2.75) {人工设计};
\node[option,fill=orange!68!black,text=white] (automated) at (8.76,-3.85) {自动答案搜索};

\draw[wire] (root.east) -- (1.46,0);
\draw[wire] (1.46,-4.20) -- (1.46,3.80);
\foreach \target/\y in {model/3.80,eng/2.70,ans/-2.00,multi/-3.10,tuning/-4.20} {
  \draw[wire] (1.46,\y) -- (\target.west);
}
\draw[wire] (eng.east) -- (4.38,2.70) |- (promptshape.west);
\draw[wire] (4.38,2.70) |- (template.west);
\draw[wire] (ans.east) -- (4.38,-2.00) |- (answershape.west);
\draw[wire] (4.38,-2.00) |- (search.west);

\draw[wire] (promptshape.east) -- (7.30,3.80) |- (cloze.west);
\draw[wire] (7.30,3.80) |- (prefix.west);
\draw[wire] (template.east) -- (7.30,1.60) |- (manual.west);
\draw[wire] (7.30,1.60) |- (auto.west);
\draw[wire] (answershape.east) -- (span.west);
\draw[wire] (7.30,-.60) |- (token.west);
\draw[wire] (7.30,-.60) |- (sentence.west);
\draw[wire] (search.east) -- (7.30,-3.30) |- (human.west);
\draw[wire] (7.30,-3.30) |- (automated.west);
\end{tikzpicture}

\caption{基于提示的学习中，各种设计选择及决策之间的关系。}
\label{fig:ch03-design}
\end{figure}

\subsection{提示形式}
提示模板大致有两种形式，即前缀提示（prefix prompt）和完形填空提示（cloze prompt）。下文用“提示形式”（prompt shape）统称这种区分。

\subsubsection{前缀提示}
在前缀提示中，输入及描述输入的文本都位于答案空位之前，形成前缀。例如，在影评情感分析任务中，前缀提示模板可以写作
\begin{equation}
\text{“}\symbf{x}'=\text{[x] This movie is [z]}\text{”},
\tag{3.5}
\end{equation}
将模板应用于输入，会得到“Cannot watch this movie. This movie is [z]”这样的提示。

如果回答与输入可以按从左到右的顺序直接对应，两者之间的联系简单明确，前缀提示往往表现良好。例如，在机器翻译任务中，明确给出目标语言，前缀提示就可以有效地引导模型生成译文。考虑以下模板。
\begin{equation}
\begin{aligned}
\text{“}\symbf{x}'={}&\text{Translate the following English sentence}\\
 &\text{to French: [x][z]}\text{”},
\end{aligned}
\tag{3.6}
\end{equation}
前缀提示通过明确指定所需语言，向模型给出直接的指令，从而产生准确的译文。

\subsubsection{完形填空提示}
与前缀提示不同，完形填空提示的答案空位两侧都有模板词元，所以答案位于模板内部。例如，在影评情感分析任务中，完形填空提示模板可以表示为
\begin{equation}
\text{“}\symbf{x}'=\text{[x] This is a [z] movie.}\text{”},
\tag{3.7}
\end{equation}
答案空位由上下文中的词语包围。把模板应用到输入上，就会得到提示“Cannot watch this movie. This is a [z] movie.”。这里的标点也算模板文本。因此，若模板在答案空位 \texttt{[z]} 之后只有一个句号，仍视为完形填空提示。

完形填空提示形式灵活，可以构造出更自然的语言结构，尤其适合多项选择问答等任务。把答案空位放在一个有上下文的短语中，可以促使模型生成符合周围语言模式的输出。答案两侧的上下文让模型能考虑输入中的语言特征和情境线索，最终生成更准确、更有意义的内容。

两种提示形式的区别如下。
\begin{itemize}
\item \textbf{前缀提示。}输入和模板文字都位于答案空位之前。例如“Cannot watch this movie. This movie is [z]”。
\item \textbf{完形填空提示。}输入和模板文字围绕着答案空位。例如“Cannot watch this movie. It was a [z] movie.”。
\end{itemize}

\subsection{手工设计模板}
现在转向模板的设计和优化。最直接的方法是手工设计模板，利用人的专业知识、直觉和对任务的理解，设计合适的模板。通常既要考虑任务本身的特点，也要用启发式方法确定合适的提示结构和措辞。为了改善表现，可能需要反复调整模板（\refcite{ch03-brown}{Brown et al., 2020}；\refcite{ch03-petroni}{Petroni et al., 2019}；\refcite{ch03-schick-2020a}{Schick and Schütze, 2020a},b）。最终应根据在有标签数据集上的表现作决定。

手工生成提示没有通用做法，最重要的原则是试验多个候选。以返回各国首都的提示为例，可以看到反复尝试的重要性。下面给出四个候选模板。
\begin{lstlisting}[caption={国家首都查询的提示模板},label={lst:ch03-capitals}]
1. "the capital city of [x] is [z] ."
2. "[z] is the capital city of [x] ."
3. "what is the capital city of [x]? It is [z] ."
4. "[z] is located in [x], and is its capital city ."
\end{lstlisting}

这四个模板看起来都说得通，效果是否一样？原书用 AllenNLP 掩码语言建模演示工具\footnote{\url{https://demo.allennlp.org/masked-lm}} 测试了一个输入样例。在这个演示工具中输入一个带掩码词元的句子，模型会返回最可能填入空位的若干词元。以 Poland 为输入 \texttt{[x]}，对每个模板预测前三个词元，结果见表~\ref{tab:ch03-capitals}。

\begin{table}[!htbp]
\centering
\caption{在 AllenNLP 掩码语言模型演示中使用清单~\ref{lst:ch03-capitals} 各模板的预测分数。给定 \texttt{[x]} 为“Poland”，每行给出填入 \texttt{[z]} 的前三个预测词元的概率，以百分数表示。}
\label{tab:ch03-capitals}
\small
\begin{tabular}{@{}lrrrrrr@{}}
\toprule
输入 & Warsaw & Kraków & Poznan & Poland & here & It\\
\midrule
模板 1 & 37.7 & 23.2 & 19.3 & — & — & —\\
模板 2 & 55.7 & 14.8 & 15.0 & — & — & —\\
模板 3 & 10.1 & — & — & 21.2 & 3.6 & —\\
模板 4 & 29.3 & — & 12.3 & — & — & 33.0\\
\bottomrule
\end{tabular}
\end{table}

模板 1 和模板 2 的首选词元是正确答案“Warsaw”，模板 2 对“Warsaw”的预测优势更大。值得注意的是，它们是四个模板中最简单直接的两个，都没有多个句子或累赘的从句安排。模板 3 的首选答案是“Poland”，模板 4 则预测代词“It”。两者的第二选择都是“Warsaw”，但显然没有像前两种模板那样有效地激发大语言模型中隐含的知识。

在手工设计提示时，设计者应当用多个模板和许多类似的有标签样例进行试验，以便作统计优化。应根据某项指标确定最优提示，例如采用选出首选预测的办法。
\begin{equation}
A(t_r,i)=
\frac{\sum_{\langle x,y\rangle\in\mathcal{R}}
 \delta\bigl(y=\arg\max_{y'}P_{\mathrm{LM}}(y'\mid x,t_{r,i})\bigr)}
{\mathcal{R}}.
\tag{3.8}
\end{equation}
这里 $\mathcal{R}$ 是由关系 $r$ 对应的带标签主语—宾语对组成的测试集。$\delta(\cdot)$ 是克罗内克的 $\delta$ 函数。当 $y$ 等于语言模型的首选预测时，函数返回 1，否则返回 0。最终选择在这组主语—宾语对训练样例上准确率最高的提示。
\transnote{原书式（3.8）的分母印为 $\mathcal{R}$。若计算准确率，分母是集合大小 $|\mathcal{R}|$。本段也保留了原书先称“测试集”、后称“训练样例”的措辞差异。}

% PDF 133--139；原书第 103--109 页。
\subsection{自动设计模板——离散搜索}
自动设计模板，在预先定义的搜索空间中用搜索或生成的方法，寻找效果最好的提示模板。这种方法实现起来更复杂，不过它一般会更充分地搜索参数空间，通常比手工提示工程表现更好。自动提示工程可分为离散搜索（discrete search）与连续搜索（continuous search）。

这些自动设计方法主要区别在于，使用离散的还是连续的词元或提示来引导语言模型。离散提示（discrete prompt）完全由自然语言词元或短语组成，连续提示（continuous prompt）则包含连续的、可调的参数。离散提示包括本节迄今所见的模板，其中连接输入 $x$ 与掩码输出 $z$ 的词元固定不变。连续提示中有不固定的词元，可在模型训练中调整。例如，可以用连续提示“[a1] [a2] [a3] [a4] [x] [a5] [z]”替换离散模板“the capital city of [x] is [z] .”，训练时微调其中的词元 $a_n$ 以改善结果。下面考察这些类别中的代表性方法及其应用前景。表~\ref{tab:ch03-discrete} 汇总了这些方法。

\subsubsection{提示挖掘}
提示挖掘（prompt mining）最早由 \refcite{ch03-jiang}{Jiang 等人（2020）}提出。它从大量文本语料中挖掘提示，依据是主语 $x$ 和宾语 $y$ 附近的词往往描述两者的关系。

\begin{keypoint}
仍以各国首都为例。在一个大型语料库中，Poland 和 Warsaw 在彼此附近出现时，通常暗示国家与首都之间的某种关系。如果收集许多具有同一种关系的主语—宾语对，例如更多国家与其首都，再抽取语料库中这些词对共同出现的句子，这些句子就可以作为生成该信息检索任务提示模板的依据。
\end{keypoint}

这种语料挖掘方法可以用两种办法生成提示。第一种称为中间词提示提取（middle-word prompt extraction）。在搜索语料中找到含有主语—宾语对的句子，抽取两者之间的文本词元，直接作为提示模板。仍以“波兰的首都是华沙”为例，目标是挖掘最能激发这条知识的提示。搜索语料库中包含这两个实体的句子，可以找到下面的例句。
\begin{quote}
\ttfamily
Warsaw is the capital city of Poland, and\\
has a population of 1.86 million people.
\end{quote}
只抽取主语—宾语对之间的词，得到“\texttt{is the capital city of}”。由此形成模板“\texttt{[z] is the capital city of [x]}”。

对小规模训练集得到的范围内所有主语—宾语对重复这个过程，搜索并抽取每一对的中间词提示。

\refcite{ch03-jiang}{Jiang 等人（2020）}的第二种挖掘方法使用更精细的语言学分析。具体而言，用句法分析从匹配的句子中抽取主语—宾语对之间最短的依存路径，将它作为模板。

例如，从句子“The capital of Poland is Warsaw.”抽取的中间词模板为“[x] is [z]”，这个模板显然太简单，无法完成任务。对同一个句子作依存分析，会得到下面的依存路径。
\[
\text{Poland}\xleftarrow{\mathrm{pobj}}\text{of}
\xleftarrow{\mathrm{prep}}\text{capital}\xleftarrow{\mathrm{nsubj}}\text{is}
\xleftarrow{\mathrm{attr}}\text{Warsaw}
\]
由此得到模板“capital of [x] is [z]”，看起来更适合查询一国首都。依存分析得到的模板也可能更容易激发目标知识，因为它们来自较为固定的语言规则，而大语言模型应已在预训练中学到这些规则（\refcite{ch03-jawahar}{Jawahar et al., 2019）}。

\subsubsection{提示改写}
提示改写（prompt paraphrasing）从已有提示出发，生成措辞尽可能多样的模板变体。仍以查询各国首都为例，可以写出若干略有区别的版本。
\begin{itemize}
\item \textbf{原始提示}\enspace “[z] is the capital city of [x] .”
\item \textbf{改写 1}\enspace “[z], the capital city of [x] .”
\item \textbf{改写 2}\enspace “[z] is the capital of [x] .”
\item \textbf{改写 3}\enspace “[x]’s capital city is [z] .”
\item \textbf{改写 $n$}\enspace “[x]’s capital, [z] .”
\end{itemize}
为自动完成这种改写，\refcite{ch03-mallinson}{Mallinson 等人（2017）}提出回译（back-translation）。先把原提示译成另一种语言，得到 $B$ 个候选译文。再把每个候选译文译回原来的语言，共得到 $B^2$ 个候选模板。按照往返翻译概率给模板排序，再从中筛选，概率计算如下。
\begin{equation}
P(t)=P_{\mathrm{forward}}(\bar{t}\mid\hat{t})\cdot P_{\mathrm{backward}}(t\mid\bar{t}),
\tag{3.9}
\end{equation}
其中 $\hat{t}$ 是原始提示，$\bar{t}$ 是原始提示 $\hat{t}$ 的译文，$t$ 是最终参与排序的候选提示。保留排名前 $T$ 的候选。之后还可进一步筛选和集成候选提示，使它们在目标自然语言处理任务中更有用。第 3.5 节会继续讨论提示筛选与集成。

\subsubsection{梯度引导搜索}
还可以用梯度引导搜索（gradient-directed search）设计提示模板。这个想法最早由 \refcite{ch03-wallace}{Wallace 等人（2019）}提出，他们当时研究对生成模型的对抗攻击（adversarial attack）。研究者设计了一种算法，把“触发词元”（trigger token）加到前缀提示末尾，紧挨答案空位之前，并迭代更新这些词元，使将错误回答填入空位时的损失最小。下面是一个问答（Question Answering，QA）提示的例子。
\begin{equation}
\text{“Question: [x] Context: [y] Answer: [T] [T] [T] }[z_{\mathrm{adv}}]\text{”},
\tag{3.10}
\end{equation}
其中 \texttt{[x]} 和 \texttt{[y]} 分别是问题与上下文，$[z_{\mathrm{adv}}]$ 是试图诱使模型产生的对抗输出，\texttt{[T]} 表示一串可变的触发词元。这些词元可迭代更新，以最小化某个语言模型对该序列的损失。更新使用基于 HotFlip 的梯度引导搜索（\refcite{ch03-ebrahimi}{Ebrahimi et al., 2018）}。这种做法诱使模型产生对抗回答。重要的是，研究者发现，在许多情形下，输入文本即使变化，优化出的触发词元序列仍然有效，会让模型针对许多不同输入生成同一个不合适的输出。下面是该研究使用 SQuAD 数据集中的一组问答所做的示例。

\begin{quote}
\textit{Question:} Why did he walk?

For exercise, Tesla walked between 8 to 10 miles per day. He squished his toes one hundred times for each foot every night, saying that it stimulated his brain cells. \textbf{why how because} \underline{to kill american people}
\end{quote}
三个词元“why how because”是梯度优化的产物。对许多不同输入，它们都会使 GPT-2 生成带下划线的对抗回答。

人们很快认识到，这种方法可以用于提示工程中的模板优化。在此基础上，\refcite{ch03-shin}{Shin 等人（2020）}提出了自动构造提示模板的方法 AutoPrompt。研究者用一组包含触发词元的初始模板，形式类似式（3.10），再在大量有标签的输入—输出样例上迭代，用梯度引导搜索优化词元。图~\ref{fig:ch03-autoprompt} 以情感分析为例说明了这一方法。送给语言模型的输入由三部分组成。
\begin{itemize}
\item \textbf{原始输入（$x_{\mathrm{inp}}$）。}对应图~\ref{fig:ch03-flow} 的输入 $x$。
\item \textbf{触发词元（$x_{\mathrm{trig}}$）。}通过梯度搜索学到的自然语言词元。学到多少词元取决于初始化梯度搜索时设定的词元数，这个数量可以视作超参数。
\item \textbf{答案空位。}在图~\ref{fig:ch03-autoprompt} 中记作 \texttt{[P]} 或 \texttt{[MASK]}，对应图~\ref{fig:ch03-flow} 示例中的 \texttt{[z]}。
\end{itemize}

\sourcefigure{fig03_11.pdf}{用 AutoPrompt 考察掩码语言模型能否进行情感分析。输入 $x_{\mathrm{inp}}$ 被放入自然语言提示 $x_{\mathrm{prompt}}$，其中有一个待预测的 \texttt{[MASK]} 词元。模板 $\lambda$ 将原始输入与学到的触发词元 $x_{\mathrm{trig}}$ 组合起来。触发词元在各个输入间共享，由上述梯度搜索过程学得。对各类别的自动选择的标签词元集合，将掩码语言模型预测的 $p(\texttt{[MASK]}\mid x_{\mathrm{prompt}})$ 求和，得到类别 $y$ 的概率。}{fig:ch03-autoprompt}

按给定提示模板的结构组合各部件，得到送入语言模型的优化输入 $x_{\mathrm{prompt}}$。将自动选择的一组输出词元的概率相加，就能确定标签类别。在本例中，正面类别对应 Cris、marvelous 和 philanthrope，负面类别对应 worse、incompetence 和 Worse。正面标签的累计概率高于负面标签，因此分类结果为正面。

\begin{keypoint}
优化得到的词元在人看来可能不直观。在同一个 LAMA T-REx 实体关系子集基准上，\refcite{ch03-shin}{Shin 等人（2020）}报告称，以手工模板为参照，AutoPrompt 比 \refcite{ch03-jiang}{Jiang 等人（2020）}评估的 Top-1 改写提示多提高了 9 个百分点的准确率。使用 BERT 和 RoBERTa 变体时，在答案映射、自然语言推断、事实检索和实体关系抽取等任务上，AutoPrompt 的表现比手工提示高 10\% 至 20\%。经过优化的提示甚至可以超过微调模型，尤其是在数据较少、只有少数有标签样例的场景中。
\end{keypoint}

\subsubsection{提示生成}
上述离散提示搜索方法都使用掩码语言模型，预测单个词元。受 \refcite{ch03-radford-2019}{Radford 等人（2019）}、\refcite{ch03-brown}{Brown 等人（2020）}及其他研究展示的上下文学习能力启发，\refcite{ch03-gao}{Gao 等人（2021）}提出了语言模型的改进少样本微调（Better Few-shot Fine-tuning of Language Models，LM-BFF）。他们用大型文本到文本模型 T5（\refcite{ch03-raffel}{Raffel et al., 2020）}取代前面例子中的掩码语言模型。其方法是以少量训练样例组成少样本提示，引导 T5 搜索合适的提示。

图~\ref{fig:ch03-lmbff} 展示 \refcite{ch03-gao}{Gao 等人（2021）}用 T5 生成提示模板的过程。先从小型训练集 $D_{\mathrm{train}}$ 抽取若干样例，送入 T5，提示它构造模板 $\mathcal{T}$，不必明确指定模板应包含多少个词元。T5 的输入经过专门设计，以引导模型生成提示。图中将 $D_{\mathrm{train}}$ 的样例按类别分组，组织成适合 T5 的输入，再生成一个模板或一组模板 $\mathcal{T}$，使下式最大。
\begin{equation}
\sum_{(x_{\mathrm{in}},y)\in D_{\mathrm{train}}}
\log P_{\mathrm{T5}}\bigl(\mathcal{T}\mid\mathcal{T}(x_{\mathrm{in}},y)\bigr),
\tag{3.11}
\end{equation}
其中 $P_{\mathrm{T5}}$ 是 T5 语言模型的输出概率分布。

\sourcefigure{fig03_12.pdf}{提示生成过程。把训练集 $D_{\mathrm{train}}$ 的输入样例按类别划分，构造为适合 T5 的输入，通过解码得到一个或少量模板，使式（3.11）的求和最大。图中先生成候选模板，再经微调和评估选出最佳模板。}{fig:ch03-lmbff}

在用这些构造出的输入使 T5 生成一组模板后，解码生成的模板。在 $D_{\mathrm{train}}$ 上对语言模型 $\mathcal{L}$ 作微调，再用开发集 $D_{\mathrm{dev}}$ 评估，选出最佳模板。\refcite{ch03-gao}{Gao 等人（2021）}表明，这种提示模板生成方法与在给定输入旁附上语义相似的示范相结合，表现显著优于手工设计的提示。此外，用 T5 配合手工生成或自动生成的提示，表现优于标准微调方案，说明了基于提示的学习在自然语言处理中的作用。

在 \refcite{ch03-gao}{Gao 等人（2021）}的工作基础上，\refcite{ch03-ben-david}{Ben-David 等人（2022）}提出即时任意领域适配的提示学习算法（Prompt learning algorithm for on-the-fly Any-Domain Adaptation，PADA）。该方法也利用 T5 生成文本到文本提示的能力，专门生成可以表示多个源领域、供人阅读的提示。针对预测分布外数据（out-of-distribution data）这一常见难题，PADA 将多个特定领域映射到同一个语义空间，以提高泛化能力。\refcite{ch03-ben-david}{Ben-David 等人（2022）}报告称，在文本分类和序列标注任务上，该方法相对于表现良好的对照方法取得了明显提升。

% 按原书第 109 页重排表 3.3，保留横置阅读方向。
\begin{table}[p]
\centering
\makebox[\linewidth][c]{\rotatebox{90}{%
\begin{minipage}{.94\textheight}
\caption{离散提示的自动设计方法。}
\label{tab:ch03-discrete}
\footnotesize
\setlength{\tabcolsep}{3pt}
\renewcommand{\arraystretch}{1.14}
\begin{tabularx}{\linewidth}{@{}>{\raggedright\arraybackslash}p{.10\linewidth}
  >{\raggedright\arraybackslash}p{.155\linewidth}
  >{\raggedright\arraybackslash}X
  >{\raggedright\arraybackslash}p{.174\linewidth}
  >{\raggedright\arraybackslash}p{.182\linewidth}@{}}
\toprule
方法 & 概述 & 示例 & 优点 & 缺点\\
\midrule
提示挖掘 &
在大型语料中搜索包含给定查询—答案对的句子，并用它们构造提示模板。 &
\textbf{测试对}\enspace Poland :: Warsaw\par
\textbf{挖掘到的句子}\enspace “Warsaw is the capital city of Poland, and has a population 1.86 million people”（见维基百科）\par
\textbf{所得模板}\enspace “[z] is the capital city of [x] .” &
程序实现简单。\par 可用更复杂的语言学抽取过程扩展。 &
所得模板受语料中已有句子的领域限制。\par 难以保证优化效果。\\
\midrule
提示改写 &
让种子提示经过多轮翻译，生成许多略有差异的变体，再选出表现最好的模板。 &
\textbf{种子提示}\enspace [z] is the capital city of [x] .\par
\textbf{回译变体}\par
“[z], the capital city of [x] .”\par
“[x]’s capital city is [z] .”\par
“[x]’s capital, [z] .” &
程序实现简单。\par 测试的提示多样，有助于优化。 &
回答覆盖范围较窄。\par 难以保证优化效果。\\
\midrule
梯度搜索 &
将一组可变的触发词元（本例为 [T]）和输入—预测对组合成模板，在训练中优化，使预测结果最好。 &
\textbf{评论／情感对}\enspace “a real joy” :: positive\par
\textbf{初始提示}\enspace “a real joy [T][T][T] [T][T] positive”\par
\textbf{梯度优化后的提示}\enspace “a real joy atmosphere alot dialogue Clone totally positive”\par
\textbf{优化后的模板}\enspace [x] atmosphere alot dialogue Clone totally [z] &
可以找到高度优化的输入词元。\par 不依赖该领域已有的句子语料。 &
计算开销大，程序实现复杂。\par 产生的模板不直观。\par 输出模板受人类语言的嵌入表示约束。\\
\midrule
提示生成 &
编码器—解码器模型（如 T5）预测种子模板中的词元。种子模板由训练查询—回答对构造。再把预测结果转成模板并评估质量。 &
\textbf{影评情感的种子模板}\par
A pleasure to watch. <X> great <Y>\par
No reason to watch. <X> terrible <Y>\par
\textbf{T5 填充后的模板}\par
A pleasure to watch. This is great .\par
No reason to watch. A terrible one.\par
\textbf{所得待测模板}\par
[x] This is [z].\par
[x] A [z] one. &
测试的提示多样，有助于优化。\par 大语言模型生成的模板本身可能更适合大语言模型使用。 &
计算开销大，程序实现复杂。\par 种子模板需要人工输入。\par 难以保证优化效果。\\
\bottomrule
\end{tabularx}
\end{minipage}}}
\end{table}


% PDF 138、140--143；原书第 108、110--113 页。
\subsection{自动设计模板——连续搜索}
构造提示的主要目标，是让大语言模型高效地完成任务。提示是否便于人理解，并非必备条件，因此无需把它限定在可供人理解的自然语言范围内（\refcite{ch03-li-2019}{Li et al., 2019）}。为此，研究者开始考察连续提示，又称软提示（soft prompt），直接在模型的嵌入空间中进行提示。

\begin{keypoint}
连续提示解决了离散提示的两项主要限制。
\begin{enumerate}
\item 模板中词的嵌入无需对应英语等自然语言中词的嵌入。
\item 模板不再局限于由预训练语言模型的参数来表示。模板可以拥有自己的参数，并使用下游任务的训练数据微调这些参数。
\end{enumerate}
\end{keypoint}

\subsubsection{前缀调优}
\refcite{ch03-li-liang}{Li 和 Liang（2021）}首次提出前缀调优（prefix tuning）。受提示的上下文学习效果启发（见第~\ref{sec:ch03-icl}~节），他们在输入文本开头添加了针对特定任务的“虚拟词元”（virtual token），见图~\ref{fig:ch03-prefix-tuning}。这些向量不代表实际词元。初始化时为它们设置适当的维数，使语言模型能像关注硬词元一样关注它们。随后，这些向量可作为连续变量参与训练，而硬词元的表示是固定的。语言模型由此可以借助前缀调优获知任务的性质，同时利用提示文本中明确的离散特征。

\refcite{ch03-li-liang}{Li 和 Liang（2021）}报告说，在数据量较少时，他们的前缀调优实验优于微调，与使用全量数据的微调相比也颇有竞争力。研究者在摘要任务中对 BART（\refcite{ch03-lewis}{Lewis et al., 2019）}做前缀调优，在表格到文本任务中对 GPT-2（\refcite{ch03-radford-2019}{Radford et al., 2019）}做同样的处理。与已有的适配器方法和全量数据微调相比，前缀调优在这两项任务上都表现良好。这些结果显示，前缀调优可用于不同类型的语言模型，实验具体展示了它在编码器—解码器模型和自回归模型上的适用性。

AutoPrompt 用训练数据集在离散空间中按梯度搜索，优化一组离散提示（见第~3.3.3.3~节）。与之类似，前缀调优用训练数据学习一组连续向量，即前缀，以最大化下式。
\begin{equation}
\max_{\phi}\log p_{\phi}(y\mid x)
  =\sum_{i\in Y_{\mathrm{idx}}}\log p_{\phi}(z_i\mid h_{<i}).
\tag{3.12}\label{eq:ch03-prefix-objective}
\end{equation}
通常，$p_{\phi}$ 由大语言模型可训练的参数 $\phi$ 决定。这里大语言模型的参数固定，因此用 $P_{\theta}$ 表示可训练的前缀参数 $\theta$。$h_i$ 是第 $i$ 个时间步各层激活值（包括前缀处激活值）的拼接。前缀调优利用从左到右生成的自回归语言模型。顾名思义，学得的向量放在模型序列最左侧，前缀的影响在模型的各个固定层中从左向右传播。

\sourcefigure{fig03_13.pdf}{上半图是预训练再微调。每做一项下游任务，都要复制并调优预训练模型的参数。有些语言模型规模很大，因此这样做开销显著，也有技术上的困难。下半图是前缀调优。它冻结模型参数，仅调优特定于任务的前缀。只需替换经过调优的前缀，同一个大语言模型便可用于多项下游任务。这样可将各任务的处理方式做成可替换的模块，并减少总体存储开销。}{fig:ch03-prefix-tuning}

\subsubsection{混合提示与用离散词元初始化的提示}
\refcite{ch03-li-liang}{Li 和 Liang（2021）}指出，前缀参数的初始化和对学习率的敏感性，可能使前缀调优的训练不稳定。他们在该项工作中将前缀参数化，以代替使用大型前馈神经网络生成的较小矩阵\footnote{译注：Li 和 Liang 论文第 4.3 节的做法是将较小的矩阵输入前馈网络来表示前缀，训练后仅保存得到的前缀。见 \url{https://aclanthology.org/2021.acl-long.353/}。}。另一种初始化连续词元的方法，是使用预先选定的离散词元。这些词元可以经由前文介绍的离散模板自动搜索学得（如 \refcite{ch03-zhong}{Zhong et al., 2021）}，也可以手工定义。在实体关系知识探查中，以这些词元为起点学习连续词元，已显示出良好前景（\refcite{ch03-qin}{Qin and Eisner, 2021}；\refcite{ch03-zhong}{Zhong et al., 2021）}。

\refcite{ch03-liu-2021}{Liu 等人（2021）}提出了另一种有前景的自动提示模板设计方法，用硬词元和软词元共同组成最终模板。图~\ref{fig:ch03-hybrid} 比较了离散提示模板（a）和这类混合提示（b）的生成过程。在 BERT 的知识探查任务中，他们的 P-tuning 方法生成的提示模板优于其他所有参与比较的离散提示模板。使用 ALBERT（\refcite{ch03-lan}{Lan et al., 2020）}时，这种方法在少样本 SuperGLUE 基准上达到了当时的最佳水平。

\sourcefigure{fig03_14.pdf}{对“The capital of Britain is [MASK]”搜索提示的例子。已知上下文（深灰色，“Britain”）和目标（白色方框，“[MASK]”），浅灰色区域表示提示词元。在（a）中，提示生成器只能收到离散奖励。在（b）中，可以采用不同方法优化伪提示和提示编码器。在某些情况下，加入少量与任务有关的锚定词元，如（b）中的“capital”，还可提高下游任务的表现。}{fig:ch03-hybrid}

\subsection{基于提示的微调}
\label{sec:ch03-prompt-finetuning}
前面介绍的提示方法，都固定用于推理的模型，通过改变提示或前缀来优化结果。另一种方法则解除对模型参数的冻结，将输入与输出对按固定模板组织成数据集，再用传统的反向传播方法微调模型参数。这种方法称为基于提示的微调（prompt-based fine-tuning）。仍以图~\ref{fig:ch03-formal-prompt} 的例子为例。
\begin{quote}
\texttt{Cannot watch this movie. This is [z].}
\end{quote}
除了可以用这个模板推理，还可以调优模型，使它准确预测人工指定的 \texttt{[z]} 值。要做好基于提示的微调，可能需要大量以同一模板组织的输入与输出对。不过，要获得良好表现，通常比常规 PTFT 所需的有标签样例少（参见第~\ref{sec:ch03-tutorial-one}~节的教程）。

这种方法既适用于完形填空式提示，也适用于前缀式提示。采用前者时，用相同形式的提示填入不同的输入与输出对来微调模型，使模型预测答案槽位中的词元。采用后者时，在答案的前缀上微调模型，使它逐步预测后续答案的下一个词元，答案可为任意长度。微调完成后，按标准的基于提示的学习流程用模型推理。

\begin{keypoint}
基于提示的微调是指令调优（instruction tuning）的基本技术，而指令调优是开发 ChatGPT 等先进聊天机器人的重要步骤。第~\ref{sec:ch04-instruction}~节将详细讨论指令调优。这里先指出，基于提示的训练一般优于零样本或少样本的基于提示的学习，代价是可能需要大量计算资源。它在许多情况下也优于 PTFT，尤其是在数据较少时。
\end{keypoint}

% PDF 143--145；原书第 113--115 页。
\section{答案工程}
\label{sec:ch03-answers}
提示工程帮助我们选择合适的模板。类似地，答案工程（answer engineering）设计和优化答案格式，引导语言模型针对具体的任务或问题，给出尽可能准确且切合上下文的回答。为此，需要考虑不同因素，确定答案形式，再用不同方法搜索答案空间，将答案映射到标签。

\subsection{答案形式}
首先要考虑答案形式（answer shape）。它决定模型输出的粒度，可以细到单个词元，也可以是完整的短语或句子。不同任务需要的答案粒度不同，选好答案形式是基于提示的学习取得成效的关键。答案形式有三种基本类型。

\begin{itemize}
\item \textbf{词元。}从预训练语言模型的词表或其子集中，取一个或多个单独的词元。这种形式常用于情感分类、关系抽取和命名实体识别等分类任务（\refcite{ch03-cui}{Cui et al., 2021}；\refcite{ch03-petroni}{Petroni et al., 2019}；\refcite{ch03-yin}{Yin et al., 2019）}。例如，在情感分类中，模型的答案可以是单个词元，如“positive”“negative”或“neutral”。采用这种形式时，答案空间通常只包含少数几个候选词元，因此属于受约束答案空间。
\item \textbf{片段。}片段（chunk 或 span）是较短的多词元序列，通常与完形填空式提示搭配使用。与词元形式不同，片段的长度不固定，一般属于无约束答案空间。这使它适合问答任务，例如回答“\texttt{Dante was born in [z] .}”这样的提示。
\item \textbf{句子。}按任务的需要，答案可以由一个或多个句子组成，也可以是整篇文档。句子形式的答案通常配合前缀式提示使用，常见于摘要、翻译等需要较详细回答的语言生成任务（\refcite{ch03-radford-2019}{Radford et al., 2019）}。这种形式的答案是无约束的。
\end{itemize}

\subsection{定义答案空间}
用 $Z$ 表示答案空间，即模型针对输入可能给出的所有答案的集合。在许多情况下，这个答案空间映射到一组输出类别，记作 $Y$。答案空间大致分为受约束和无约束两类。

\begin{keypoint}
\begin{itemize}
\item 在受约束答案空间中，输出限于一组预定义的答案。对于输出标签数目有限的任务，如文本分类、实体识别或选择题问答，这种空间很有用。
  \begin{itemize}
  \item 在这种设置中，$Z$ 中的每个元素都映射到标签空间 $Y$ 中的一个元素，作为最终输出。例如，在图~\ref{fig:ch03-mask} 的情感分析任务中，答案“terrible”映射到负面类别。
  \item 这种空间与词元形式的答案相关。
  \end{itemize}
\item 在无约束答案空间中，$Z$ 包含所有词元（\refcite{ch03-petroni}{Petroni et al., 2019）}、定长片段（\refcite{ch03-jiang}{Jiang et al., 2020）}或词元序列（\refcite{ch03-radford-2019}{Radford et al., 2019）}。
  \begin{itemize}
  \item 一般无需把这些输出映射到单独的标签空间，它们本身就是最终输出。
  \item 这种空间与片段及句子形式的答案相关。
  \end{itemize}
\end{itemize}
\end{keypoint}

如上所述，受约束答案空间需要选出组成 $Z$ 的具体词元。以下三节讨论怎样选这些元素。

\subsection{手工建立答案映射}
定义 $Z$ 最简单的办法是手工选择。这需要反复试验。先为各个标签选出不同的答案集合，再用数据检验它们的效果。以情感分析的例子为例，提示为“$x'=$ \texttt{[x] This is a [z] movie.}”。最初取
$Z=\{\text{great}\to\text{positive},\ \text{terrible}\to\text{negative}\}$。
也可以试验
$Z=\{\text{good}\to\text{positive},\ \text{bad}\to\text{negative}\}$。
哪一种表现较好，须由实验确定。

映射不必是一一对应的。多个答案可以对应同一个标签。例如，可以把答案空间扩展为 $Z=\{\text{great},\ \text{terrible},\ \text{good},\ \text{bad}\}$，看概率最高的答案属于正面类（great 或 good），还是负面类（terrible 或 bad），据此选择标签。

\subsection{自动建立答案映射——离散搜索}
有研究者开发了自动选择答案的技术。\refcite{ch03-jiang}{Jiang 等人（2020）}提出一种迭代方法，从初始答案空间 $Z'$ 出发，通过改写扩大它的范围。研究者收集答案—标签对 $(z',y)$，对答案 $z'$ 作回译，找出意义相近、措辞不同的词，再检验这些词的效果。

\refcite{ch03-gao}{Gao 等人（2021）}也提出了一种定义标签词映射 $Z\to Y$ 的方法，目标是使小型验证集上的准确率最大。将填入输入数据的模板交给大语言模型，由模型给出可能的映射，再汇集似然最高的预测。这样，数据集的每个标签都有一份词元排序表，其中的词元是大语言模型预测可用于填补掩码的。可以取前 $n$ 项来建立答案映射。

\subsection{自动建立答案映射——连续搜索}
连续答案搜索直接在模型的嵌入空间中进行。和连续提示一样（见第~3.3.4~节），它解除输出必须映射到已知词元的限制，为进一步优化留出空间。以“positive”和“negative”为例。在自然语言中，这两个词或许最接近所需的输出，但它们的嵌入未必是模型解决该任务时最合适的输出。连续搜索允许模型按更贴近其内部信息表示的方式作答，无需将答案强制转换为自然语言。

% PDF 146--151；原书第 116--121 页。
\section{多提示推理}
\label{sec:ch03-multi}
本章迄今讨论的提示形式都有一个隐含的共同点，即由一个查询和一处待填的回答组成，回答可以包含一个或多个词元。例如，“\texttt{The capital city of Poland is [MASK]}”只提出一个问题，也只要求一个答案。这样推理时，模型只能利用其预训练权重和这一条提示。若上下文十分有限，大语言模型的预测能力可能受到限制。在零样本场景下，模型尚未针对所请求的任务微调，这种限制尤为明显。此外，即使模板已经过优化，特定措辞带来的偏差仍可能造成预测不准确或难以消除的系统性错误。

为解决这些问题，已有文献探讨了多种多提示推理（multi-prompt inference）方法。这些方法会在推理时把问答对作为额外上下文交给模型，或汇总不同提示或子提示的结果，以改善总体表现。下面简要介绍其中几种，并用实例说明它们的优点。

\subsection{集成}
集成（ensembling）是最简单的多提示推理方法。对一个给定的查询—答案对 $(x,z)$，使用多个提示模板 $f_{\mathrm{prompt}}$，再合并各次预测的分数，得到该推理问题的总体最佳猜测。不同的提示模板各有细微的偏差，单独使用时可能猜错，汇总后则可能比任一个模板都更准确。

为说明集成的效果，我们重新使用第~3.3.2~节的任务和提示，仍用 AllenNLP 查询加拿大首都。四条掩码语言提示如下。
\begin{lstlisting}[caption={查询加拿大首都的提示},label={lst:ch03-canada}]
1. [MASK] is located in Canada, and is its capital city.
2. I am in the capital of Canada, I am in [MASK].
3. [MASK] is the capital city of Canada.
4. The capital city of Canada is [MASK].
\end{lstlisting}

表~\ref{tab:ch03-ensemble} 列出每条提示分数最高的五项，其中包括几个加拿大城市，也包括“Canada”本身。提示 1 正确预测 Ottawa（渥太华），不过只略高于第二名 Winnipeg（温尼伯）。提示 3 情况相近。提示 4 识别首都 Ottawa 的效果最好。提示 2 将 Ottawa 排在第四，在多个其他加拿大城市之后。最下一行给出各个不同的回答词元在这些提示间的平均预测分数，未进入前五名时，按零分计（见表注）。集成结果正确给出了加拿大首都，其前两名的分数差大于提示 1、2 和 3。

\begin{table}[!htbp]
\centering
\caption{按清单~\ref{lst:ch03-canada} 的提示，用 AllenNLP 进行零样本推理时，若干加拿大城市及“Canada”的预测分数。最下一行是各列的简单平均值。未进入前五名的城市以“—”表示，求平均时计作 0。由于只记录了前五项，这些平均值并非严格的全量平均值。例如，提示 2 很可能对 Winnipeg 给出非零分数，计入后该城市的平均分会略增。这个影响很小，不改变结论。}
\label{tab:ch03-ensemble}
\small
\setlength{\tabcolsep}{2.5pt}
\begin{tabular}{@{}lrrrrrrrrr@{}}
\toprule
输入 & Otta. & Winn. & Mont. & Toro. & Calg. & Lond. & Canada & Edmo. & Vanc.\\
\midrule
提示 1 & 17.3 & 15.8 & 10 & 6.8 & 8.1 & — & — & — & —\\
提示 2 & 7.3 & — & 5.5 & 11.3 & — & 13.4 & 12.2 & — & —\\
提示 3 & 24.9 & 19.8 & 10 & 10.4 & 7.2 & — & — & — & —\\
提示 4 & 29.8 & 10.6 & 14.4 & — & — & — & — & 7.6 & 6.2\\
平均 & 19.83 & 11.55 & 9.98 & 7.13 & 3.83 & 3.35 & 3.05 & 1.9 & 1.55\\
\bottomrule
\end{tabular}
\end{table}

用于集成的一组模板可通过多种途径得到，第~3.3.2--3.3.4~节介绍了其中几种。除了选择哪些提示参与集成，还要考虑怎样把多个提示算出的概率合成一个数值。文献中考察了几种分数汇总方法。

\begin{itemize}
\item \textbf{简单多数投票。}选择在最多提示中排在第一的回答 \texttt{[z]}（如 \refcite{ch03-lester}{Lester et al., 2021）}。长期以来，研究表明简单多数投票在许多场景下的表现与更复杂的集成方法相当（\refcite{ch03-lam}{Lam and Suen, 1997}），可以用它来比较更复杂的集成结果。
\item \textbf{简单平均。}对每个不同的回答 \texttt{[z]}，把集成中所有填好的提示模板 $x'$ 上的概率 $P(z\mid x')$ 求和，再除以提示总数。然后选取平均分最高的回答。
\item \textbf{加权平均。}可以为集成中的不同提示赋予权重，增大或减小它们在平均值中的贡献。这可以看作简单平均的推广。权重可按对表现的主观判断临时设定，也可根据测试集预测准确率等客观指标选择（如 \refcite{ch03-schick-2020a}{Schick and Schütze, 2020a}），还可与基于提示的学习一起调优（\refcite{ch03-jiang}{Jiang et al., 2020）}。
\item \textbf{知识蒸馏。}最后一种方法是先在无标签数据集上生成集成预测，把这些预测用作另一个模型的训练输入，再以这个模型作为预测任务的分类器。推理时将集成缩减为一个模型，可以降低开销。由于新模型学习了集成结果的特征，还能保有集成的好处（\refcite{ch03-hinton}{Hinton et al., 2015）}。\refcite{ch03-schick-2021}{Schick 和 Schütze（2021）}等人已将这种方法用于提示学习。
\end{itemize}

\subsection{上下文学习}
\label{sec:ch03-icl}
多提示的第二种方法是上下文学习（in-context learning，ICL），又称提示扩充（prompt augmentation）或示例学习（demonstration learning）（\refcite{ch03-liu-2023}{Liu et al., 2023）}。在给定提示之前或之后附上一个或多个填好查询—答案对的模板实例，再把得到的多提示字符串交给模型推理。例如，可以把“\texttt{The capital of Canada is [MASK] .}”扩充为“\texttt{The capital of France is Paris. The capital of Canada is [MASK] .}”。利用语言模型识别模式的能力，增加上下文能大幅提高大语言模型在分类和生成任务上的准确率（如 \refcite{ch03-brown}{Brown et al., 2020）}。

以下是 AllenNLP 对上一段的两条查询给出的分数。

\noindent\begin{minipage}{\linewidth}
\begin{lstlisting}[caption={以加拿大和法国首都为例的上下文学习},label={lst:ch03-icl}]
The capital of Canada is [MASK] .
"Ottawa": 0.28438693284988403
"Montreal": 0.1832585334777832
...

The capital of France is Paris. The capital of Canada is [MASK] .
"Ottawa": 0.3411746919155121
"Vancouver": 0.1890593022108078
...
\end{lstlisting}
\end{minipage}\par

填入一个上下文示例后，排在第一的正确答案分数略有增加，说明上下文学习有所帮助。选择哪些样例来扩充上下文，对准确程度的提升影响很大。用 \refcite{ch03-lu}{Lu 等人（2022）}的话说，这种选择“可使表现从近乎当时最佳水平跌至随机猜测”。例如，把另一个上下文样例放在待答提示的前面或后面，会得到以下结果。
\begin{lstlisting}[caption={扩展上下文示例，包含法国、泰国和加拿大首都},label={lst:ch03-icl-extended}]
The capital of France is Paris. The capital of Thailand is Bangkok. The capital of Canada is [MASK] .
"Vancouver": 0.31520649790763855
"Ottawa": 0.277070015668869
...

The capital of France is Paris. The capital of Canada is [MASK] . The capital of Thailand is Bangkok.
"Ottawa": 0.4077857434749603
"Vancouver": 0.2586808502674103
...
\end{lstlisting}

把新增的上下文样例放在待答提示前面时，模型误把 Vancouver（温哥华）认作加拿大首都。若把它放在末尾，模型表现则进一步改善。预测分数对上下文样例及其位置的细微变化十分敏感。因此，构造上下文学习的提示模板时需要仔细考虑，也可以考虑自动优化（如 \refcite{ch03-lu}{Lu et al., 2022）}。

\subsection{提示分解}
最后一种多提示学习方法，是将一个复杂的提示任务拆成多个独立的子任务，为每个子任务分别构造提示。如果输入一串文本需要得到多个回答，这种方法通常很适用。例如，命名实体识别要给输入句子中不止一个词元分配标签。更一般地说，这种方法适用于词性识别等序列标注任务。

以命名实体识别为例，看下面这个句子。
\begin{lstlisting}[caption={Serena 在法拉盛草地获胜},label={lst:ch03-serena}]
Serena won the tennis tournament at Flushing Meadows.
\end{lstlisting}

若想识别这句话中的每个命名实体，就很难用一条提示同时指出 Serena 是人名、tennis 是一种运动、Flushing Meadows 是地点。完形填空式提示尤其难做到这一点。可以把任务简化为一组提示，每条提示询问句子中的一个词元。

\begin{lstlisting}[caption={提示分解示例},label={lst:ch03-decomp}]
Serena won the tennis tournament at Flushing Meadows. Choosing between name, location , sport, or none, "Serena" is a [MASK] .
"nickname": 0.17498865723609924
"choice": 0.042668603360652924

Serena won the tennis tournament at Flushing Meadows. Choosing between name, location , sport, or none, tennis is a [MASK] .
"sport": 0.08051449060440063,
"synonym": 0.03813215345144272

Serena won the tennis tournament at Flushing Meadows. Choosing between name, location , sport, or none, "Flushing Meadows" is a [MASK] .
"suburb": 0.1842394471168518,
"neighborhood": 0.09695208817720413,

Serena won the tennis tournament at Flushing Meadows. Choosing between name, location , sport, or none, "the" is a [MASK] .
"suffix": 0.11233088374137878
"noun": 0.08212842047214508
\end{lstlisting}

仅用这种粗略的方法，模型就能较准确地识别三个实体，并对非实体词元“the”给出非实体回答。在实际的命名实体识别系统中，可为输入句子的每个 $n$ 元组构造一条提示并生成回答，否则会漏掉 Flushing Meadows 这类由多个词元组成的实体。配合答案映射器（verbalizer）限制允许的选项，并用一系列同样拆成子提示的输入样例调优模型，可以进一步提高准确率。

表~\ref{tab:ch03-multi} 汇总了本节介绍的三种方法及其优缺点。到这里，我们已经介绍了基于提示的学习文献中几种重要的提示工程方法。接下来通过实际操作，看看如何在自己的项目中试用这些方法。

% 按原书第 121 页重排表 3.5。
\begin{table}[p]
\centering
\makebox[\linewidth][c]{\rotatebox{90}{%
\begin{minipage}{.94\textheight}
\caption{多提示推理方法。}
\label{tab:ch03-multi}
\footnotesize
\setlength{\tabcolsep}{3pt}
\renewcommand{\arraystretch}{1.14}
\begin{tabularx}{\linewidth}{@{}>{\raggedright\arraybackslash}p{.095\linewidth}
  >{\raggedright\arraybackslash}p{.173\linewidth}
  >{\raggedright\arraybackslash}X
  >{\raggedright\arraybackslash}p{.16\linewidth}
  >{\raggedright\arraybackslash}p{.17\linewidth}@{}}
\toprule
方法 & 概述 & 示例 & 优点 & 缺点\\
\midrule
集成 &
用多个模板完成一项推理任务，以简单多数投票、平均、蒸馏等办法，将每个测试样例的各次预测汇总成最终答案。 &
提示 1\enspace “I am in the capital of Canada, I am in [MASK] .”\par
提示 2\enspace “[MASK] is located in Canada, and is its capital city .”\par
提示 3\enspace “[MASK] is the capital of Canada”\par
各提示的首选结果\enspace 1) Toronto，2) Ottawa，3) Ottawa\par
简单多数投票\enspace Ottawa（正确） &
对多个提示求平均，可以减轻个别提示失效带来的影响。\par
可充分利用提示工程方法生成的多个候选模板。 &
推理时，每个测试样例需要多次预测。\par
效果受输入提示的质量限制，优化余地较小。\\
\midrule
上下文学习 &
推理前，在尚未填入答案的模板之前或之后添加已填好查询—答案对的模板示例，引导大语言模型给出所需的预测。 &
原提示\enspace “The capital of Canada is [MASK] .”\par
上下文提示\enspace “The capital of France is Paris. The capital of Canada is [MASK] .”\par
两种提示的首选结果\enspace 1) Ottawa（$p=0.284$），2) Ottawa（$p=0.341$） &
推理时，每个测试样例只需预测一次。\par
实现简单、直观。 &
准确率很大程度上取决于所选的上下文示例。\par
要得到好的效果，可能需要对上下文提示微调。\\
\midrule
提示分解 &
构造多个子提示，分别探查复杂查询的不同方面，例如识别一句话中的各个命名实体。 &
原提示\enspace \texttt{[x]} = “Serena won the tennis tournament”\par
分解后的提示\par
“[x] . ‘Serena’ is a [MASK]”（人名）\par
“[x] . ‘Tennis’ is a [MASK]”（运动）\par
“[x] . ‘the’ is a [MASK]”（非实体） &
可探查查询中粒度更细的细节。 &
推理时，每个测试样例需要多次预测。\par
可能需要逐一优化多个子提示。\\
\bottomrule
\end{tabularx}
\end{minipage}}}
\end{table}


% PDF 150、152--160；原书第 120、122--130 页。
\section{教程一——比较文本分类与命名实体识别中的提示和预训练再微调方法}
\label{sec:ch03-tutorial-one}
\subsection{概述}
本章介绍了基于提示的学习的概念，并详细讨论了提示与答案形式的几种配置。与 PTFT 相比，基于提示的方法的主要优点之一，是它在零样本和少样本条件下的表现。这一点还没有演示过。本教程将说明，与传统的基于任务头的微调相比，基于提示的学习怎样用更少的训练样例取得更好的结果。借助这一特点，可以用较少的数据和算力让大语言模型适应新任务。

\begin{keypoint}
\noindent\textbf{目标}
\begin{itemize}
\item 比较基于提示的学习与基于任务头的微调。
\item 演示怎样组织提示以完成不同的任务。
\item 介绍 OpenPrompt 库，说明本章介绍的技术是怎样实现的。
\item 绘制学习曲线，展示提示在少样本设置下的表现。
\end{itemize}
\end{keypoint}

这里给出教程的简明版本。完整版见 \url{https://github.com/springer-llms-deep-dive/llms-deep-dive-tutorials}。

\subsection{实验设计}
实验将直接比较 PTFT 和基于提示的学习在零样本及少样本条件下做文本分类和命名实体识别的能力。两种微调实验都以 BERT 为基础，使用 PyTorch，基于提示的部分还使用 OpenPrompt。逐步增加训练数据子集的大小，依次调优 BERT 模型，并定期在验证集上做预测，观察模型在少样本条件下怎样学习。最后，比较两种调优方法在各项自然语言处理任务上的学习曲线，并讨论其含义。

\subsection{结果与分析}
\subsubsection{文本分类}
首先，用传统的基于任务头的方法微调预训练语言模型，做文本分类。如第~3.1.2~节所述，这一过程用句子—标签对调优任务专用预测头，将语言模型（这里是 BERT）丰富的语言表示迁移到分类任务。数据集是 GLUE 中用于情感分析的 SST-2，由从电影评论中抽取的句子组成。每句话都由人工标注，表明它表达的是对电影的正面情感还是负面情感。把训练集分成正负两类，使用不同数量的样例训练模型，再在测试集上评估表现。

结果见表~\ref{tab:ch03-sst2} 的中列。零样本准确率接近 50\%，相当于随机猜测。在最初的几个样本规模上，准确率提高不多，到 64 个样例时有明显改善，最后在 256 个样例时达到 87\%。表中数值对实验具体选取哪些训练样例很敏感，多次运行所得的数值会略有不同，但总体趋势不变。

接着换用基于提示的训练进行比较。所得准确率见表~\ref{tab:ch03-sst2} 右列。每类仅用 32 个样例，提示模型的准确率就大幅上升。此后曲线很快趋于平缓，增加数据带来的收益较小。

\begin{table}[!htbp]
\centering
\caption{在 GLUE SST-2 数据集上，比较预训练再微调和基于提示的文本分类。表中给出不同训练样本数下的预测准确率。}
\label{tab:ch03-sst2}
\begin{tabular}{@{}rrr@{}}
\toprule
训练样本数 & PTFT 准确率 & 提示方法准确率\\
\midrule
0 & 0.5092 & 0.6800\\
16 & 0.5069 & 0.6743\\
32 & 0.6548 & 0.7867\\
64 & 0.8486 & 0.8475\\
128 & 0.8624 & 0.8521\\
256 & 0.8739 & 0.8658\\
\bottomrule
\end{tabular}
\end{table}

\sourcefigure{fig03_15.pdf}{表~\ref{tab:ch03-sst2} 的数据图。图中比较文本分类实验里预训练再微调与基于提示的学习的准确率随训练样例数的变化。}{fig:ch03-sst2}

图~\ref{fig:ch03-sst2} 从零样本出发，逐渐增加训练数据，绘出两种模型的准确率随训练样例数的变化。与基于任务头的分类器相比，提示模型只用极少的训练样例就取得了出色的结果。PTFT 模型最终可与基于提示的模型竞争，但要每类使用 32 个训练样例，才能达到与提示模型零样本表现相近的水平。\footnote{译注：原文称 32 个样例时已达到提示模型的零样本准确率。表~\ref{tab:ch03-sst2} 所列的两个数值分别为 0.6548 和 0.6800。}

\subsubsection{命名实体识别}
现在考察第二项分类任务，即命名实体识别。预计它在零样本条件下表现较差。我们要确定，在少样本条件下，提示方法是否仍优于微调。实验使用 CoNLL-2003 数据集，它收集了新闻标题和引语，逐词元标记人名、组织名、地名，另设一类容纳其他实体。

PTFT 的结果见表~\ref{tab:ch03-conll} 中列。与文本分类实验类似，BERT 在没有任务训练时表现较差。这并不令人意外，因为命名实体识别要识别标签，而 BERT 还不知道该数据集中各标签所指的含义。它预测的值大多大于 0，而多数标签等于 0，因而产生许多假阳性，$F_1$ 得分也较低。最初几个数据点显示，加入更多训练样例时，表现先变差，可能源于灾难性遗忘。随后，每增加一批数据，表现便有所改善，到 1024 个样例时达到 $F_1=85\%$。

\begin{table}[!htbp]
\centering
\caption{在 CoNLL-2003 数据集上，比较预训练再微调与基于提示的命名实体识别。表中给出不同训练样本数下的 $F_1$ 得分。}
\label{tab:ch03-conll}
\begin{tabular}{@{}rrr@{}}
\toprule
训练样本数 & PTFT 的 $F_1$ & 提示方法的 $F_1$\\
\midrule
0 & 0.0687 & 0.0712\\
8 & 0.0492 & 0.5788\\
16 & 0.0023 & 0.6482\\
32 & 0.0034 & 0.7274\\
64 & 0.3323 & 0.7867\\
128 & 0.5578 & 0.8365\\
256 & 0.7157 & 0.8672\\
512 & 0.7894 & 0.8304\\
1024 & 0.8526 & 0.8551\\
\bottomrule
\end{tabular}
\end{table}

接着采用基于提示的调优，逐次增加训练样例数，每次训练五轮，并查看学习曲线。表~\ref{tab:ch03-conll} 右列给出结果。仅用几句话，表现就大幅改善，并逐步提高，到 256 个样例时 $F_1$ 达到 87\%。使用更多数据调优也未能使表现进一步提高。

图~\ref{fig:ch03-conll} 汇总比较了 PTFT 和基于提示的调优。结果与文本分类类似。数据充足时，预训练再微调可与基于提示的学习竞争。数据缺乏时，基于提示的调优则取得了好得多的结果。

\sourcefigure{fig03_16.pdf}{表~\ref{tab:ch03-conll} 的数据图。图中比较命名实体识别实验里预训练再微调和基于提示的学习的 $F_1$ 得分随训练样例数的变化。}{fig:ch03-conll}

\subsection{结论}
本教程的两个实验共同表明，训练集规模有限时，基于提示的调优优于基于任务头的微调。取得规模足够大、质量良好的训练数据，长期以来都是许多机器学习问题的主要难点。因此，少样本学习取得的效果尤其值得注意。采用提示方法时，模型的质量取决于提示模板和答案格式的设计。选对这两项，只需几十个样例就能得到良好的自然语言处理结果。下一章将进一步讨论这一点。

\section{教程二——提示工程方法}
\label{sec:ch03-tutorial-two}
\subsection{概述}
本章的另一项主题是用模板工程提高经过提示调优的大语言模型的分析能力。清单~\ref{lst:ch03-capitals} 和表~\ref{tab:ch03-capitals} 说明，大语言模型的推理输出对模板结构及提示组成的细节很敏感。该演示使用了一个网页应用，适合验证思路，但因为不能微调，能力有限。本教程将通过少样本和多样本提示调优扩展这些实验，讨论不同提示模板的效果，说明优化提示模板的重要性。

\begin{keypoint}
\noindent\textbf{目标}
\begin{itemize}
\item 说明任务表现对提示模板设计十分敏感，即便细微改动，也能造成明显差异。
\item 考察哪些因素有助于提高提示模板的质量。
\item 用软提示进行自动调优，比较它与手工提示的效果。
\end{itemize}
\end{keypoint}

这里给出教程的简明版本。完整版见 \url{https://github.com/springer-llms-deep-dive/llms-deep-dive-tutorials}。

\subsection{实验设计}
本教程考察几种不同的模板工程方法，使用基准数据集训练模型并评估其表现。先用最简单的手工模板工程方法（见第~3.3.2~节）。这种方法由开发者构造最适合任务的模板。可以参考针对不同提示学习任务推荐模板的文献（例如第~3.2.4~节），也可以试验不同配置。

我们还会用训练数据探索自动模板设计。第~3.3.3 和 3.3.4~节已经介绍了几种自动方法。本教程采用一种基于梯度的模板优化方法，称为软提示。手工提示固定，软提示则使用可变模板，可以在不受离散词元映射限制的情况下将其调到更合适的表示。先用一个组合了数据集特征与“软”词元的模板初始化软提示，其中软词元可选择用特定的词或短语初始化。然后用反向传播调整相应嵌入，以达到训练目标。我们用 OpenPrompt 的代码实现软提示（\refcite{ch03-ding}{Ding et al., 2021）}。实验将说明，软提示可以优于手工构造的提示。

本教程用 SuperGLUE BoolQ 数据集。它提供由一段说明性文字、一条与段落内容有关的是非题和正确答案组成的三元组。BoolQ 的主题范围很广，包括历史、科学、地理、法律、体育和流行文化等，适合考察大语言模型的自然语言推断能力。

\subsection{结果与分析}
在手工提示实验中，我们用三种不同的提示模板和若干种训练样本数，对 \texttt{t5-base} 模型做基于提示的调优。训练样例在两个标签类别之间均分。我们把三种模板分别称为“最简单”模板、“最简单加标点”模板和更“合适”的模板。第一种只是把段落、问题和掩码连接起来。第二种用标点进一步指明文本结构。第三种加入文字提示各部分的含义，并把掩码词元放在形式合理的完形填空中。

\begin{lstlisting}[language=Python,nolol=true]
temp1 = "{passage} {question} {mask}"
temp2 = "{passage} . {question} ? {mask}"
temp3 = "hypothesis: {passage} premise: {question} The answer was {mask} ."
\end{lstlisting}

首先用最简单的模板做微调实验，结果见表~\ref{tab:ch03-boolq} 左侧的准确率列。总体来看，模型表现较差。

\begin{itemize}
\item 零样本推理将所有样本都判为负类，因此准确率等于验证集的负类样本比例（约为 50\%）。
\item 少样本结果有所改善，但只略好于随机猜测，相当于每道题抛一枚硬币。
\item 模型表现大约在 256 个样例时达到最高，准确率仍然较低。
\end{itemize}

\begin{table}[!htbp]
\centering
\caption{三种提示模板在 SuperGLUE BoolQ 上的零样本和少样本准确率。}
\label{tab:ch03-boolq}
\begin{tabular}{@{}rrrr@{}}
\toprule
训练样本数 & 最简单 & 最简单加标点 & 改进模板\\
\midrule
0 & 0.501 & 0.526 & 0.564\\
16 & 0.511 & 0.508 & 0.489\\
32 & 0.513 & 0.482 & 0.503\\
64 & 0.518 & 0.516 & 0.535\\
128 & 0.526 & 0.598 & 0.638\\
256 & 0.578 & 0.628 & 0.686\\
512 & 0.546 & 0.638 & 0.640\\
\bottomrule
\end{tabular}
\end{table}

接着试一项简单的改动，在段落末尾没有句号时加上句号，并在问题结尾加上问号。表~\ref{tab:ch03-boolq} 中间的准确率列显示了这个小改动的效果。只需在恰当的位置加入“?”和“.”，零样本准确率便由约 50\% 增至约 53\%。微调开始后，带标点模板的准确率上升得比最简单的模板快，说明提示有所改善。

最后试用更合适的手工模板。这个模板给出了有用的上下文，并明确向模型提问，预期效果更好。用改进后的模板重复相同实验，得到表~\ref{tab:ch03-boolq} 右列。图~\ref{fig:ch03-boolq} 给出三条学习曲线。改进模板的零样本表现明显优于另两种模板。不过，微调样例很少时，灾难性遗忘使准确率下降，继续微调之后才开始回升。使用 256 个样例后，模型能正确回答约 69\% 的提示，明显优于另两种模板。

\sourcefigure{fig03_17.pdf}{使用三种不同提示训练的模型在 SuperGLUE BoolQ 验证集上的准确率变化。模板形式会影响问答任务的零样本和少样本表现。}{fig:ch03-boolq}

接着用精心设计的 10 个候选模板重复上述实验，其中既有完形填空式，也有前缀式（详见完整版教程）。结果显示，模板有几个值得注意的特点。同一类模板的表现有一定共性。

\begin{itemize}
\item 前缀式提示在零样本条件下已有一定效果。用于调优的样例很少时，灾难性遗忘使表现下降。样本数进一步增加后，预测能力大幅提高。
\item 完形填空式提示的零样本表现稍差，只有少量训练样例时表现也会变差。随后，它的表现会改善，最终与前缀式提示相当。
\item 提供的上下文较少的提示，表现明显较差。零样本结果只略好于随机猜测，使用全部训练样例后，也只再提高几个百分点。不过，完全训练后，每条提示的表现都优于随机猜测，说明模型确实编码了其中一部分问题的答案。
\item 同一类中的不同模板，总体表现有明显差异，且随训练数据增加，差异往往也会增大。这说明模板结构中的细微差别也可能带来显著影响。
\end{itemize}

最后一个实验构造了两种软提示，一个采用简单模板，一个采用经过设计的模板。两种情况下，都用 128 个 SuperGLUE BoolQ 样例微调若干轮。这个实验冻结 \texttt{t5-base} 模型，只调优提示。图~\ref{fig:ch03-soft-boolq} 左图显示两种模型在验证集上的表现变化。

\begin{itemize}
\item 红色虚线代表缺少特征的提示。微调 60 轮后，表现仍无提高。这个模板的内容过少，缺少可以微调的特征，因此微调未能改善效果。
\item 黑色实线代表第二种提示。继续微调后，表现明显提高，60 轮内准确率约提高 4.5\%。与输入模板相比，软提示调优得到了更好的模板。
\end{itemize}

经过 60 轮训练，经过设计的软提示模型优于手工模板实验中所有模型的零样本表现，也优于这些模型用 128 个训练样本微调之前的表现。由于提示的规模远小于大语言模型本身，达到这一表现所需的算力较少。因此，当数据点有限、需要考虑训练算力开销时，提示调优可能优于大语言模型调优。

接着考察四种调优方法。

\begin{enumerate}
\item 仅调优软提示。
\item 仅调优大语言模型。
\item 先调优软提示，再调优大语言模型。
\item 同时调优软提示和大语言模型。
\end{enumerate}

图~\ref{fig:ch03-soft-boolq} 右图给出结果\footnote{译注：原图右侧横轴标为训练轮数，原文在下段按训练样例数讨论这些点。}，可以看到以下现象。

\begin{itemize}
\item 对于 BoolQ 任务，调优大语言模型明显优于单独调优软提示。只调提示的模型在继续调优后确实有明显改善，但在 512 个样例时，其表现远低于另外三个允许调整大语言模型参数的模型。随着训练时间延长，这一差距应会缩小。\refcite{ch03-lester}{Lester 等人（2021）}采用经过优化的训练方案，经过 30\,000 步得到约 0.9 的准确率。
\item 其他三个模型中，以调优过的软提示初始化的模型，零样本表现最好。它之前已经接触过训练数据，因此这一点并不出人意料。此后，允许调整大语言模型的几个模型，随训练样本增加而改善的速度相近。各个训练阶段，同时调优大语言模型和提示的模型略优于只调优大语言模型的模型。这很可能是因为两者共同调优时可调整的参数更多（2.48 亿对 2.22 亿），它也是本教程中表现最好的模型。不过，与使用经过设计的提示并调优大语言模型的做法相比，提升幅度很小。
\end{itemize}

\sourcefigure{fig03_18.pdf}{左图是从简单提示和经过设计的提示出发，进行软提示调优的结果。右图是图例所列四种学习方式的学习曲线。同时调优提示和大语言模型的模型，在训练过程的各个阶段表现最好。}{fig:ch03-soft-boolq}

\subsection{结论}
实验表明，提示工程对改善大语言模型表现很重要。要取得最佳表现，当然还有许多可调因素，这里没有重点考察，例如训练集规模、训练轮数、学习率和大语言模型的选择等。尽管如此，从表现最差的模型到表现最好的模型，仅通过改变模板设计，预测准确率就提高了 25\% 以上。因此，使用提示的模型需要认真考虑模板设计。

% PDF 160--163；原书第 130--133 页。
\begin{chapterreferences}
\refentry{ch03-ben-david}{Eyal Ben-David, Nadav Oved, and Roi Reichart. PADA: Example-based prompt learning for on-the-fly adaptation to unseen domains, 2022.}
\refentry{ch03-brown}{Tom Brown, Benjamin Mann, Nick Ryder, Melanie Subbiah, Jared D Kaplan, Prafulla Dhariwal, Arvind Neelakantan, Pranav Shyam, Girish Sastry, Amanda Askell, et al. Language models are few-shot learners. \textit{Advances in Neural Information Processing Systems}, 33:1877–1901, 2020.}
\refentry{ch03-cui}{Leyang Cui, Yu Wu, Jian Liu, Sen Yang, and Yue Zhang. Template-based named entity recognition using BART. \textit{arXiv preprint arXiv:2106.01760}, 2021.}
\refentry{ch03-ding}{Ning Ding, Shengding Hu, Weilin Zhao, Yulin Chen, Zhiyuan Liu, Hai-Tao Zheng, and Maosong Sun. OpenPrompt: An open-source framework for prompt-learning, 2021.}
\refentry{ch03-ebrahimi}{Javid Ebrahimi, Anyi Rao, Daniel Lowd, and Dejing Dou. HotFlip: White-box adversarial examples for text classification, 2018.}
\refentry{ch03-gao}{Tianyu Gao, Adam Fisch, and Danqi Chen. Making pre-trained language models better few-shot learners. In \textit{Proceedings of the 59th Annual Meeting of the Association for Computational Linguistics and the 11th International Joint Conference on Natural Language Processing (Volume 1: Long Papers)}, pages 3816–3830, Online, August 2021. Association for Computational Linguistics. doi: 10.18653/v1/2021.acl-long.295. URL \url{https://aclanthology.org/2021.acl-long.295}.}
\refentry{ch03-hinton}{Geoffrey Hinton, Oriol Vinyals, and Jeff Dean. Distilling the knowledge in a neural network, 2015.}
\refentry{ch03-jawahar}{Ganesh Jawahar, Benoît Sagot, and Djamé Seddah. What does BERT learn about the structure of language? In \textit{Proceedings of the 57th Annual Meeting of the Association for Computational Linguistics}, pages 3651–3657, Florence, Italy, July 2019. Association for Computational Linguistics. doi: 10.18653/v1/P19-1356. URL \url{https://aclanthology.org/P19-1356}.}
\refentry{ch03-jiang}{Zhengbao Jiang, Frank F Xu, Jun Araki, and Graham Neubig. How can we know what language models know? \textit{Transactions of the Association for Computational Linguistics}, 8:423–438, 2020.}
\refentry{ch03-kamath}{Uday Kamath, John Liu, and James Whitaker. Transfer learning: Domain adaptation. \textit{Deep Learning for NLP and Speech Recognition}, pages 495–535, 2019.}
\refentry{ch03-kotsiantis}{Sotiris B Kotsiantis, Ioannis Zaharakis, P Pintelas, et al. Supervised machine learning: A review of classification techniques. \textit{Emerging Artificial Intelligence Applications in Computer Engineering}, 160(1):3–24, 2007.}
\refentry{ch03-kumar}{Ankit Kumar, Ozan Irsoy, Peter Ondruska, Mohit Iyyer, James Bradbury, Ishaan Gulrajani, Victor Zhong, Romain Paulus, and Richard Socher. Ask me anything: Dynamic memory networks for natural language processing. In Maria Florina Balcan and Kilian Q. Weinberger, editors, \textit{Proceedings of The 33rd International Conference on Machine Learning}, volume 48 of Proceedings of Machine Learning Research, pages 1378–1387, New York, New York, USA, 20–22 Jun 2016. PMLR. URL \url{https://proceedings.mlr.press/v48/kumar16.html}.}
\refentry{ch03-lam}{L. Lam and S. Y. Suen. Application of majority voting to pattern recognition: An analysis of its behavior and performance. \textit{Trans. Sys. Man Cyber. Part A}, 27(5):553–568, sep 1997. ISSN 1083-4427. doi: 10.1109/3468.618255. URL \url{https://doi.org/10.1109/3468.618255}.}
\refentry{ch03-lan}{Zhenzhong Lan, Mingda Chen, Sebastian Goodman, Kevin Gimpel, Piyush Sharma, and Radu Soricut. ALBERT: A lite BERT for self-supervised learning of language representations, 2020.}
\refentry{ch03-lester}{Brian Lester, Rami Al-Rfou, and Noah Constant. The power of scale for parameter-efficient prompt tuning. \textit{arXiv preprint arXiv:2104.08691}, 2021.}
\refentry{ch03-lewis}{Mike Lewis, Yinhan Liu, Naman Goyal, Marjan Ghazvininejad, Abdelrahman Mohamed, Omer Levy, Ves Stoyanov, and Luke Zettlemoyer. BART: Denoising sequence-to-sequence pre-training for natural language generation, translation, and comprehension, 2019.}
\refentry{ch03-li-liang}{Xiang Lisa Li and Percy Liang. Prefix-tuning: Optimizing continuous prompts for generation. \textit{arXiv preprint arXiv:2101.00190}, 2021.}
\refentry{ch03-li-2019}{Xiaoya Li, Jingrong Feng, Yuxian Meng, Qinghong Han, Fei Wu, and Jiwei Li. A unified MRC framework for named entity recognition. \textit{arXiv preprint arXiv:1910.11476}, 2019.}
\refentry{ch03-liu-2023}{Pengfei Liu, Weizhe Yuan, Jinlan Fu, Zhengbao Jiang, Hiroaki Hayashi, and Graham Neubig. Pre-train, prompt, and predict: A systematic survey of prompting methods in natural language processing. \textit{ACM Computing Surveys}, 55(9):1–35, 2023.}
\refentry{ch03-liu-2021}{Xiao Liu, Yanan Zheng, Zhengxiao Du, Ming Ding, Yujie Qian, Zhilin Yang, and Jie Tang. GPT understands, too, 2021.}
\refentry{ch03-lu}{Yao Lu, Max Bartolo, Alastair Moore, Sebastian Riedel, and Pontus Stenetorp. Fantastically ordered prompts and where to find them: Overcoming few-shot prompt order sensitivity, 2022.}
\refentry{ch03-mallinson}{Jonathan Mallinson, Rico Sennrich, and Mirella Lapata. Paraphrasing revisited with neural machine translation. In \textit{Proceedings of the 15th Conference of the European Chapter of the Association for Computational Linguistics: Volume 1, Long Papers}, pages 881–893, Valencia, Spain, April 2017. Association for Computational Linguistics. URL \url{https://aclanthology.org/E17-1083}.}
\refentry{ch03-och}{Franz Josef Och, Daniel Gildea, Sanjeev Khudanpur, Anoop Sarkar, Kenji Yamada, Alexander Fraser, Shankar Kumar, Libin Shen, David A Smith, Katherine Eng, et al. A smorgasbord of features for statistical machine translation. In \textit{Proceedings of the Human Language Technology Conference of the North American Chapter of the Association for Computational Linguistics: HLT-NAACL 2004}, pages 161–168, 2004.}
\refentry{ch03-peters}{Matthew E Peters, Sebastian Ruder, and Noah A Smith. To tune or not to tune? Adapting pretrained representations to diverse tasks. \textit{arXiv preprint arXiv:1903.05987}, 2019.}
\refentry{ch03-petroni}{Fabio Petroni, Tim Rocktäschel, Patrick Lewis, Anton Bakhtin, Yuxiang Wu, Alexander H Miller, and Sebastian Riedel. Language models as knowledge bases? \textit{arXiv preprint arXiv:1909.01066}, 2019.}
\refentry{ch03-qin}{Guanghui Qin and Jason Eisner. Learning how to ask: Querying LMs with mixtures of soft prompts. In \textit{Proceedings of the 2021 Conference of the North American Chapter of the Association for Computational Linguistics: Human Language Technologies}, pages 5203–5212, Online, June 2021. Association for Computational Linguistics. doi: 10.18653/v1/2021.naacl-main.410. URL \url{https://aclanthology.org/2021.naacl-main.410}.}
\refentry{ch03-radford-2018}{Alec Radford, Karthik Narasimhan, Tim Salimans, Ilya Sutskever, et al. Improving language understanding by generative pre-training. 2018.}
\refentry{ch03-radford-2019}{Alec Radford, Jeffrey Wu, Rewon Child, David Luan, Dario Amodei, Ilya Sutskever, et al. Language models are unsupervised multitask learners. \textit{OpenAI blog}, 1(8):9, 2019.}
\refentry{ch03-raffel}{Colin Raffel, Noam Shazeer, Adam Roberts, Katherine Lee, Sharan Narang, Michael Matena, Yanqi Zhou, Wei Li, and Peter J. Liu. Exploring the limits of transfer learning with a unified text-to-text transformer, 2020.}
\refentry{ch03-schick-2020a}{Timo Schick and Hinrich Schütze. Exploiting cloze questions for few shot text classification and natural language inference. \textit{arXiv preprint arXiv:2001.07676}, 2020a.}
\refentry{ch03-schick-2020b}{Timo Schick and Hinrich Schütze. It’s not just size that matters: Small language models are also few-shot learners. \textit{arXiv preprint arXiv:2009.07118}, 2020b.}
\refentry{ch03-schick-2021}{Timo Schick and Hinrich Schütze. Few-shot text generation with pattern-exploiting training, 2021.}
\refentry{ch03-shin}{Taylor Shin, Yasaman Razeghi, Robert L. Logan IV, Eric Wallace, and Sameer Singh. AutoPrompt: Eliciting Knowledge from Language Models with Automatically Generated Prompts. In \textit{Proceedings of the 2020 Conference on Empirical Methods in Natural Language Processing (EMNLP)}, pages 4222–4235, Online, November 2020. Association for Computational Linguistics. doi: 10.18653/v1/2020.emnlp-main.346. URL \url{https://aclanthology.org/2020.emnlp-main.346}.}
\refentry{ch03-wallace}{Eric Wallace, Shi Feng, Nikhil Kandpal, Matt Gardner, and Sameer Singh. Universal adversarial triggers for attacking and analyzing NLP. In \textit{Proceedings of the 2019 Conference on Empirical Methods in Natural Language Processing and the 9th International Joint Conference on Natural Language Processing (EMNLP-IJCNLP)}, pages 2153–2162, Hong Kong, China, November 2019. Association for Computational Linguistics. doi: 10.18653/v1/D19-1221. URL \url{https://aclanthology.org/D19-1221}.}
\refentry{ch03-xiong}{Caiming Xiong, Stephen Merity, and Richard Socher. Dynamic memory networks for visual and textual question answering. In \textit{International Conference on Machine Learning}, pages 2397–2406. PMLR, 2016.}
\refentry{ch03-yin}{Wenpeng Yin, Jamaal Hay, and Dan Roth. Benchmarking zero-shot text classification: Datasets, evaluation and entailment approach. \textit{arXiv preprint arXiv:1909.00161}, 2019.}
\refentry{ch03-zhang}{Yue Zhang and Joakim Nivre. Transition-based dependency parsing with rich non-local features. In \textit{Proceedings of the 49th Annual Meeting of the Association for Computational Linguistics: Human Language Technologies}, pages 188–193, 2011.}
\refentry{ch03-zhong}{Zexuan Zhong, Dan Friedman, and Danqi Chen. Factual probing is [mask]: Learning vs. learning to recall, 2021.}
\end{chapterreferences}

